% ============================================================
% Evaluación y comparación de plugins CNI para Kubernetes
% Holman Alba · Duwan Sierra
% Ingeniería Telemática — formato IEEE journal (two-column)
% ============================================================
\documentclass[journal]{IEEEtran}

% --- Encoding & idioma ---
\usepackage[utf8]{inputenc}
\usepackage[T1]{fontenc}
\usepackage[spanish,es-noindentfirst]{babel}

% --- URLs ---
\usepackage{url}

% --- Gráficos ---
\usepackage{graphicx}
\graphicspath{{./fig/}}

% --- Matemáticas ---
\usepackage{amsmath}
\usepackage{amssymb}

% --- Tablas ---
\usepackage{array}
\usepackage{tabularx}
\usepackage{booktabs}
\usepackage{colortbl}
\newcolumntype{Y}{>{\raggedright\arraybackslash}X}

% --- Código fuente ---
\usepackage{listings}
\usepackage{xcolor}
\lstset{
  basicstyle=\ttfamily\scriptsize,
  breaklines=true,
  columns=fullflexible,
  frame=single,
  backgroundcolor=\color{gray!8}
}

% --- Citas IEEE ---
\usepackage{cite}

% --- Hiperenlaces (cargar al final del preámbulo) ---
\usepackage[hidelinks,
            pdftitle={Evaluacion y comparacion de plugins CNI para Kubernetes},
            pdfauthor={Holman Alba, Duwan Sierra}]{hyperref}

% Separación silábica en español
\hyphenation{Ku-ber-ne-tes mi-cro-ser-vi-cios te-le-má-ti-ca}

% ============================================================
\begin{document}

\title{Evaluación y comparación de plugins CNI para Kubernetes:
       un enfoque orientado a QoS, seguridad y eficiencia de
       recursos en entornos de microservicios}

\author{Holman~Alba~y~Duwan~Sierra%
\thanks{Programa de Ingeniería Telemática.
Trabajo de grado de pregrado.}%
\thanks{Manuscript: \today.}}

\markboth{Ingeniería Telemática}%
{Alba \& Sierra: Evaluación y comparación de plugins CNI para Kubernetes}

\maketitle

% --- Abstract ---
\begin{abstract}
Kubernetes no impone un plugin de red específico, lo que obliga a los
equipos a seleccionar un \emph{Container Network Interface} (CNI) sin
métricas comparativas objetivas. Esta situación conduce a
sobredimensionamiento de infraestructura (CPU/RAM) y a posturas de
seguridad incompletas. El presente trabajo construye evidencia empírica
reproducible —infraestructura como código, benchmarks automatizados y
modelo de decisión multicriterio (MCDA)— para guiar esa selección.
Se evalúan Flannel, Calico, Cilium y Antrea en un clúster K3s en alta
disponibilidad sobre DigitalOcean, midiendo \emph{throughput} TCP,
latencia de establecimiento, retransmisiones y overhead de CPU/RAM del
agente CNI, así como el impacto de políticas de red Zero Trust.
Los resultados alimentan un modelo MCDA parametrizable por perfil de
uso (URLLC, Edge/IoT, microservicios transaccionales) y un prototipo
web que expone las recomendaciones.
\end{abstract}

\begin{IEEEkeywords}
Kubernetes, CNI, QoS, iperf3, Calico, Cilium, Flannel, Antrea,
NetworkPolicy, Zero Trust, MCDA, K3s, DigitalOcean, GitOps, eBPF.
\end{IEEEkeywords}

\IEEEpeerreviewmaketitle

% ============================================================
% CAPÍTULOS
% ============================================================
\section{Marco Teórico}

\subsection{El modelo de red de Kubernetes}

Kubernetes no implementa por sí mismo la conectividad entre cargas de trabajo: define un modelo de red y delega su realización en componentes externos. Comprender ese modelo es condición previa para interpretar las métricas de este trabajo, ya que el plano de red ---y en particular el plugin que lo implementa--- determina el rendimiento, la seguridad y el consumo de recursos del clúster \cite{ref2}.

El modelo de red de Kubernetes establece tres propiedades fundamentales: cada Pod recibe una dirección IP propia, los Pods se comunican entre sí sin traducción de direcciones, y los agentes de cada nodo alcanzan a todos los Pods del clúster. Sobre estas garantías se construyen las abstracciones de descubrimiento, balanceo y control de acceso descritas a continuación \cite{ref2}.

\subsubsection{Identidad de red del Pod}

Kubernetes adopta el modelo \emph{IP-por-Pod}: cada Pod obtiene su propia dirección IP y se comunica directamente con otros Pods sin traductores intermedios. Dado que la identidad del Pod es efímera ---las réplicas se crean, destruyen y sustituyen de manera dinámica---, la plataforma requiere abstracciones estables para el descubrimiento y el direccionamiento de los destinos \cite{ref2}.

\subsubsection{Nodos y reenvío de paquetes}

Cada nodo hospeda múltiples Pods y aporta CPU, memoria y conectividad. En este nivel se instalan las rutas y los mecanismos de reenvío que garantizan la entrega de paquetes al Pod de destino, conforme al modelo de red del clúster: alcance entre Pods, ausencia de superposición de rangos de direcciones y conectividad \emph{east--west} \cite{ref2}. El agente del plugin de red se ejecuta precisamente en este nivel; por ello, su consumo de CPU y memoria por nodo Worker constituye una de las métricas centrales de la comparativa.

\subsubsection{Direccionamiento estable: Services y EndpointSlices}

Un Service expone un identificador estable ---por ejemplo, una \emph{ClusterIP} o dirección IP virtual (VIP)--- que representa a un conjunto de Pods equivalentes. El cliente se conecta a ese identificador y la plataforma balancea la solicitud hacia una de las réplicas disponibles, desacoplando a los consumidores de los cambios en el ciclo de vida de los Pods \cite{ref3,ref13}. Para mantener actualizada la relación de destinos, Kubernetes emplea \emph{EndpointSlices}, que enumeran las direcciones y puertos de los Pods que respaldan cada Service y reducen la sobrecarga de control frente a los \emph{Endpoints} tradicionales \cite{ref1}. Esta capa es relevante para interpretar la latencia medida: el tiempo en que una solicitud alcanza el Pod correcto comprende tanto la resolución de la VIP como el reenvío realizado por el plugin de red.

\subsubsection{Ingreso de tráfico externo}

El tráfico procedente de Internet no accede directamente a los Pods, sino a través de una capa de entrada. La \emph{Gateway API} formaliza los roles asociados ---infraestructura, operador y equipo de aplicación--- y las reglas de capa 7 ---host, rutas y TLS--- que determinan el encaminamiento de cada solicitud hacia el Service correspondiente \cite{ref17}. En numerosos despliegues, esta capa se encuentra precedida por un balanceador de carga externo de capa 3/4; diseños de alto rendimiento como \emph{Maglev} muestran cómo el \emph{hash} consistente y el seguimiento de flujos permiten una distribución estable y resiliente antes de que el tráfico alcance el clúster \cite{ref5}.

\subsubsection{Aislamiento lógico: Namespaces y NetworkPolicies}

Un \emph{Namespace} es una partición lógica del clúster que organiza y nombra recursos por equipo o entorno (\texttt{dev}, \texttt{test}, \texttt{prod}) y evita colisiones de nombres sin multiplicar clústeres; no constituye, sin embargo, una barrera de tráfico \cite{ref2}. La unicidad de nombres opera dentro del namespace, por lo que pueden coexistir dos servicios homónimos en namespaces distintos; el nombre DNS completo de un servicio incluye el namespace (\texttt{<servicio>.<namespace>.svc.cluster.local}), y su uso evita ambigüedades de enrutamiento \cite{ref13,ref2}.

Por defecto, la conectividad \emph{east--west} entre namespaces permanece abierta. La restricción del tráfico se especifica mediante \emph{NetworkPolicies} ---definidas con \emph{labels} y puertos---, cuya aplicación efectiva depende de que el plugin de red las implemente en su plano de datos. Este mecanismo materializa el principio de \emph{Zero Trust}: conceder únicamente los flujos necesarios y verificar explícitamente cada acceso \cite{ref4,ref24,ref25,ref26}.

\subsubsection{Recorrido de una solicitud}

El trayecto de una solicitud externa integra las abstracciones anteriores:

\begin{enumerate}
\item El cliente externo alcanza la capa de entrada (Ingress/Gateway), opcionalmente a través de un balanceador de carga previo \cite{ref17}.
\item La capa de entrada aplica reglas de capa 7 y dirige la solicitud al Service correspondiente \cite{ref3,ref13}.
\item El Service consulta sus EndpointSlices y selecciona un Pod activo \cite{ref1}.
\item El plugin de red encamina el tráfico hasta el Pod de destino en su nodo \cite{ref2,ref4,ref18}.
\end{enumerate}

En síntesis, Kubernetes define identidad (IP-por-Pod), direccionamiento estable (Services y EndpointSlices), organización (Namespaces), control de acceso (NetworkPolicies) e ingreso gobernado (Gateway API); no obstante, la conectividad subyacente que sostiene todas estas abstracciones la provee el plugin de red, objeto de estudio de este trabajo \cite{ref1,ref2,ref3,ref4,ref13,ref17,ref18}.

\subsection{Container Network Interface (CNI)}

\subsubsection{Especificación y funciones núcleo}

El \emph{Container Network Interface} (CNI) es la especificación ---y el conjunto de \emph{plugins} que la implementan--- que habilita la conectividad de red de los Pods en Kubernetes. La especificación define cómo el \emph{runtime} de contenedores solicita a un plugin que conecte un Pod a la red y cómo libera los recursos al eliminarlo, mediante operaciones encadenables (\texttt{ADD} y \texttt{DEL}) que permiten combinar responsabilidades, por ejemplo un plugin para IPAM y otro para el mapeo de puertos \cite{ref18,ref4,ref21,ref20}. Sus funciones núcleo son tres:

\begin{itemize}
  \item \textbf{Identidad de red del Pod.} Asignar a cada Pod su dirección IP propia y garantizar la alcanzabilidad \emph{east--west} \cite{ref2}.
  \item \textbf{Conexión y enrutamiento.} Crear interfaces virtuales, instalar rutas y coordinar con el \emph{host} el tránsito del tráfico hacia y desde cada Pod \cite{ref18,ref4}.
  \item \textbf{Gestión de direcciones (IPAM).} Reservar y entregar direcciones desde rangos predefinidos, evitando colisiones \cite{ref20,ref18}.
\end{itemize}

La elección del CNI y de su \emph{dataplane} condiciona la latencia, el \emph{throughput} y el consumo de CPU del clúster, con efecto directo sobre el sobredimensionamiento de recursos. Por tanto, su selección informada constituye un problema de ingeniería con consecuencias económicas y de seguridad medibles \cite{ref7,ref8,ref11,ref22}.

\subsubsection{Modelos de dataplane: overlay y ruteo nativo}

El \emph{dataplane} de un CNI sigue habitualmente uno de dos enfoques. En el modo \emph{overlay}, el tráfico entre nodos se encapsula ---por ejemplo, mediante VXLAN sobre UDP--- para transportar paquetes de Pods sin modificar la red física; ofrece simplicidad operativa y desacopla el direccionamiento del clúster respecto del \emph{underlay}, a costa de cabeceras adicionales, ajustes de MTU y cierta carga de CPU por encapsulación y descapsulación. En el modo de ruteo nativo, los Pods son alcanzables mediante rutas de capa 3 entre nodos, sin túneles; suele reducir la latencia y el consumo de CPU, pero exige una integración explícita con la red subyacente y un plan de direccionamiento coherente. La implementación concreta del \emph{fast path} ---conmutación basada en OVS o programas eBPF--- condiciona adicionalmente el rendimiento y la observabilidad del plano de datos \cite{ref7,ref8,ref11,ref22}.

En este trabajo, los cuatro CNI evaluados operan en modo \emph{overlay} sobre la VPC de DigitalOcean ---Flannel, Calico y Cilium mediante VXLAN, y Antrea mediante GENEVE---, lo que acota las diferencias de encapsulación al dataplane de cada plugin y no a la topología subyacente. La evaluación se restringe al modo \emph{overlay} por ser el único compatible con la VPC de DigitalOcean sin modificar el \emph{underlay}; la comparación con ruteo nativo se plantea como línea de trabajo futuro.

\subsubsection{Gestión de direcciones (IPAM)}

La identidad de red del Pod se materializa en una dirección IP asignada por el mecanismo de \emph{IP Address Management} (IPAM) del CNI. Las estrategias de asignación pueden ser por nodo ---cada nodo administra un rango local--- o globales ---un \emph{pool} común---, con distintas implicaciones en simplicidad operativa y aprovechamiento de direcciones. El plugin \emph{host-local} ilustra el enfoque por nodo: asigna direcciones desde ficheros locales, garantizando unicidad sin dependencias externas. Un plan de direccionamiento bien definido ---rangos, \emph{leases}, reutilización y limpieza--- minimiza colisiones, facilita \emph{dual-stack} y reduce incidencias durante el escalado o la rotación de Pods \cite{ref18,ref20}.

\subsubsection{MTU y fragmentación en redes con encapsulación}

La encapsulación propia de los \emph{overlays} reduce la MTU efectiva, ya que parte de cada paquete se destina a cabeceras adicionales. Si la MTU no se ajusta, pueden aparecer fragmentación, caídas de \emph{throughput} y aumentos de latencia. La práctica recomendada consiste en calcular la MTU efectiva restando el \emph{overhead} del túnel y validar el ajuste mediante pruebas controladas, procedimiento especialmente relevante en despliegues con VXLAN o cifrado a nivel de red \cite{ref12}.

\subsubsection{Observabilidad del plano de datos}

La operación y optimización de la red del clúster requiere visibilidad sobre los flujos de tráfico: origen, destino, puertos y volúmenes. Las soluciones de \emph{flow visibility} a nivel CNI recolectan y agregan registros de flujo y los correlacionan con identidades lógicas ---namespaces y \emph{labels}---, lo que facilita detectar rutas críticas, validar políticas y priorizar mejoras. En arquitecturas de microservicios, esta observabilidad de red complementa las métricas de aplicación para diagnosticar dependencias inesperadas \cite{ref23,ref6}.

\subsubsection{Métricas de rendimiento y sobredimensionamiento}

La evaluación de un CNI debe considerar la latencia ---promedio y colas P95/P99---, el \emph{throughput} sostenido y el costo de CPU por conexión o paquete, fijando el modo del dataplane y el perfil de tráfico para garantizar comparaciones reproducibles. La literatura reporta patrones consistentes entre modos de dataplane y entre \emph{fast paths}, con variaciones significativas en entornos de recursos limitados o de alta tasa de conexión \cite{ref7,ref8,ref9,ref11,ref22,ref27}.

Estas métricas tienen una consecuencia económica directa: un CNI adecuadamente seleccionado y sintonizado reduce los recursos aprovisionados por precaución. Tres prácticas resultan especialmente efectivas: seleccionar el dataplane según el entorno, ajustar la MTU para evitar fragmentación, y automatizar \emph{NetworkPolicies} bajo el principio de \emph{Zero Trust} para limitar el tráfico innecesario. La observabilidad de flujos cierra el ciclo de mejora continua y permite sostener los objetivos de nivel de servicio (SLO) con menor consumo de CPU y memoria \cite{ref7,ref8,ref12,ref23,ref24,ref25}.

\section{Diseño de la Solución}

Este capítulo describe, de manera detallada y comprensible, cómo se diseñó e implementó la solución propuesta para comparar las tecnologías de red que se utilizan dentro de un clúster Kubernetes. Para quienes no están familiarizados con el tema, un clúster Kubernetes es como una ciudad de servidores donde muchas aplicaciones conviven y se comunican entre sí; la ``red'' de esa ciudad es lo que garantiza que cada aplicación pueda hablar con las demás de forma rápida, segura y eficiente. Al finalizar este capítulo, el lector comprenderá por qué cada decisión de diseño fue necesaria para garantizar que los resultados de la comparativa sean válidos, reproducibles y aplicables en contextos reales.

El diseño presentado en este capítulo responde directamente a las causas del problema identificadas en el árbol de problemas: la falta de criterios objetivos para seleccionar el componente de red (CNI), la ausencia de visibilidad sobre el tráfico que circula dentro del clúster, y el sobredimensionamiento de recursos que resulta de operar sin datos reales. La solución se construyó sobre cuatro pilares fundamentales: (1) una arquitectura de red bien definida, (2) un sistema de medición de Calidad de Servicio (QoS), (3) un modelo matemático para la toma de decisiones y (4) una estrategia de seguridad automatizada. Esta correspondencia directa entre causas del problema y decisiones de diseño garantiza que cada componente del sistema tenga una justificación técnica derivada del problema de investigación, y no sea un elemento arbitrario de la solución \cite{ref29}.

\subsection{Arquitectura de Red y Topología}

Antes de poder medir o comparar tecnologías, es necesario construir un entorno de pruebas que sea justo para todas ellas. Imagine que quiere saber cuál de varios automóviles consume menos gasolina: para hacer una comparación honesta, todos deben correr en la misma pista, con el mismo tipo de carretera y las mismas condiciones climáticas. En redes, ocurre algo muy similar: el entorno donde se realizan las pruebas debe eliminar cualquier factor externo que pueda alterar los resultados.

\subsubsection{Diseño de la Red Underlay (La Infraestructura Física)}

La red underlay es la base sobre la que todo lo demás se construye; es, en términos simples, la ``carretera'' real por donde viajan los datos. Para este proyecto se tomó una decisión de diseño fundamental: utilizar la nube pública de DigitalOcean en lugar de servidores físicos propios. Esta decisión tiene una razón técnica muy clara: eliminar el ``ruido'' que generaría el hardware local. En un laboratorio físico, factores como la temperatura, el estado del cable de red o interferencias electromagnéticas pueden afectar las mediciones. Al usar servidores en la nube, todos estos factores quedan controlados por el proveedor y son equivalentes para todas las pruebas \cite{ref2}. Esta equivalencia de condiciones es el requisito fundamental para que cualquier diferencia medida en throughput, latencia o consumo de CPU pueda atribuirse exclusivamente al CNI evaluado, y no a variaciones del entorno físico.

El aprovisionamiento de esta infraestructura se realizó mediante Terraform, una herramienta de Infraestructura como Código (IaC). Esto significa que toda la configuración de los servidores está escrita en archivos de texto, como una receta de cocina, lo que garantiza que el entorno se pueda reproducir exactamente de la misma manera en cualquier momento. Esta reproducibilidad es una condición necesaria para la validez científica del experimento: si los resultados no pueden recrearse bajo las mismas condiciones, la comparación entre CNIs pierde su valor probatorio \cite{ref7}.

\begin{quote}
\textbf{¿Qué es K3s en Alta Disponibilidad?}

K3s es una versión liviana de Kubernetes, ideal para entornos donde los recursos computacionales son limitados o para laboratorios de prueba. ``Alta Disponibilidad'' (HA) significa que el sistema está diseñado para seguir funcionando incluso si uno de sus servidores principales falla. En este proyecto, se configuró 1 nodo Master con base de datos PostgreSQL externa gestionada y 2 nodos Worker dedicados. Esta configuración garantiza que las pruebas de red se realicen en un ambiente estable y representativo de lo que usaría una empresa real.
\end{quote}

La topología de red diseñada incluye una red privada VPC (Virtual Private Cloud) para el tráfico interno entre nodos. Esta decisión es crítica desde el punto de vista telemático: al usar una red privada, se garantiza que la latencia medida durante las pruebas dependa exclusivamente de la eficiencia del CNI bajo evaluación y no de la variabilidad de internet público \cite{ref8}.

\begin{table*}[htbp]
\centering
\scriptsize
\begin{tabularx}{\textwidth}{|Y|Y|}
\hline
\textbf{Componente} & \textbf{Descripción y Justificación} \\
\hline
Proveedor Cloud & DigitalOcean --- entorno reproducible y controlado, elimina ruido de hardware local \\
Aprovisionamiento & Terraform (IaC) --- configuración declarativa, reproducible y versionada en Git \\
Distribución K8s & K3s --- versión liviana de Kubernetes, ideal para laboratorio de comparativas \\
Plano de control & 1 Nodo Master con datastore PostgreSQL externo --- elimina SPOF del plano de control \\
Nodos de trabajo & 2 Workers dedicados --- separación clara de roles, mediciones sin interferencia \\
Red interna & VPC privada --- aísla el tráfico inter-nodo, garantiza mediciones limpias de QoS \\
\hline
\end{tabularx}
\caption{Componentes del diseño de la red underlay y su justificación técnica.}
\end{table*}

\subsubsection{Diseño de la Red Overlay (Los CNI bajo Evaluación)}

Si la red underlay es la carretera física, la red overlay es el sistema de señalización, semáforos y rutas que organiza el tráfico dentro del clúster. Esta capa es implementada por los plugins CNI. Un CNI es el componente que hace posible que dos aplicaciones dentro del clúster se puedan ``hablar'' entre sí, aunque estén en servidores diferentes. Kubernetes no incluye esta funcionalidad por defecto; en cambio, permite que las organizaciones elijan el CNI que mejor se adapte a sus necesidades \cite{ref4,ref18}. Esta libertad de elección es precisamente el problema central que este proyecto aborda: sin criterios objetivos basados en mediciones controladas, esa decisión se toma con base en preferencias informales o recomendaciones genéricas, con consecuencias directas sobre el rendimiento y la seguridad del clúster.

El núcleo del diseño de este proyecto es la comparativa de cuatro tecnologías de encapsulamiento y ruteo, cada una con una filosofía diferente:

\begin{table*}[htbp]
\centering
\scriptsize
\begin{tabularx}{\textwidth}{|Y|Y|Y|Y|}
\hline
\textbf{CNI} & \textbf{Tecnología Base} & \textbf{Versión} & \textbf{Explicación} \\
\hline
Flannel & VXLAN (L2) & nativo K3s & CNI más sencillo. Funciona como un túnel: envuelve los datos en un ``sobre'' adicional. Fácil de instalar pero sin soporte nativo de NetworkPolicies. \\
Calico & VXLAN (DO VPC) & 3.29.1 Tigera Operator & Tigera Operator con \texttt{encapsulation: VXLAN} para DigitalOcean VPC. Políticas avanzadas L3/L4. BGP disponible en bare-metal. \\
Cilium & eBPF + VXLAN tunnel & 1.16.5 & Modo \texttt{routingMode=tunnel}, \texttt{tunnelProtocol=vxlan}. Opera directamente en el kernel con programas eBPF. Filtrado L7 HTTP nativo. \\
Antrea & OVS (SDN) & 2.2.0 & Open vSwitch con sistema jerárquico de Tiers (Emergency $>$ SecurityOps $>$ NetworkOps $>$ Platform $>$ Application). \\
\hline
\end{tabularx}
\caption{Los cuatro CNI evaluados, su tecnología base, versión desplegada y enfoque.}
\end{table*}

\textbf{VXLAN (Virtual Extensible LAN):} Es un protocolo que permite crear redes virtuales sobre redes físicas existentes. Funciona encapsulando los paquetes de datos originales dentro de paquetes UDP, lo que permite que tráfico de red de nivel 2 viaje sobre redes de nivel 3. El costo es un overhead de 50 bytes por paquete, lo que puede reducir el throughput efectivo cuando se transmiten muchos datos pequeños \cite{ref12}.

\textbf{eBPF (Extended Berkeley Packet Filter):} Es una tecnología que permite ejecutar programas directamente dentro del kernel de Linux de manera segura y eficiente. En el contexto de redes, esto significa que Cilium puede interceptar, analizar y redirigir paquetes de red sin necesidad de copiarlos al espacio de usuario, eliminando una de las principales fuentes de latencia de los sistemas tradicionales \cite{ref11}.

\textbf{OVS (Open vSwitch):} Es un switch virtual sofisticado orientado a SDN \cite{ref22}. Antrea lo utiliza para implementar conectividad y políticas con un sistema de Tiers jerárquico que permite reglas corporativas de mayor prioridad que las NetworkPolicies estándar de Kubernetes. La diversidad de enfoques representada por estos cuatro CNIs ---desde la encapsulación básica de Flannel hasta el filtrado en kernel de Cilium--- garantiza que los resultados del experimento sean representativos del espectro de opciones disponibles para un equipo DevOps real, cubriendo desde simplicidad operativa hasta capacidades avanzadas de observabilidad y seguridad.

\subsection{Diseño del Sistema de Medición de Calidad de Servicio (QoS)}

En telemática, la Calidad de Servicio (QoS) es el conjunto de métricas que determinan qué tan bien funciona una red para las aplicaciones que la utilizan. El sistema de medición diseñado en este proyecto es el corazón experimental de la investigación. Sin mediciones precisas y reproducibles, cualquier comparación entre CNIs sería subjetiva e inútil para la toma de decisiones técnicas. El diseño sigue los principios de la metrología de redes: control de variables, reproducibilidad y representatividad estadística \cite{ref6}.

\subsubsection{Metodología de Inyección de Tráfico}

Para generar tráfico de red de manera controlada, se utilizó iperf3, la herramienta de referencia en el ámbito académico y profesional para pruebas de rendimiento de redes. iperf3 funciona con un modelo cliente-servidor: un Pod emisor (cliente) envía tráfico a un Pod receptor (servidor), y ambos registran estadísticas detalladas de la transferencia.

\begin{quote}
\textbf{Solución al problema de co-localización: podAntiAffinity}

Para garantizar que el emisor y el receptor siempre estén en servidores diferentes, se utilizó una funcionalidad de Kubernetes llamada \texttt{podAntiAffinity} (anti-afinidad de pods). Se configuró como \texttt{preferredDuringSchedulingIgnoredDuringExecution} (anti-afinidad suave, weight 100), lo que significa que el scheduler \emph{prefiere} ubicar cliente y servidor en nodos distintos. Con 2 workers dedicados disponibles, el scheduler los separa en la práctica en todos los casos. Se usa \texttt{preferred} en lugar de \texttt{required} para mantener resiliencia ante reinicios de nodo.
\end{quote}

Además del aislamiento, se diseñó un esquema de automatización mediante CronJobs de Kubernetes con ráfagas de tráfico TCP de 300 segundos cada 30 minutos. Esta frecuencia permite capturar el comportamiento de la red en diferentes momentos del día, obteniendo así un perfil estadístico completo \cite{ref7}.

\subsubsection{Variables Telemáticas Capturadas}

\begin{table*}[htbp]
\centering
\scriptsize
\begin{tabularx}{\textwidth}{|Y|Y|Y|}
\hline
\textbf{Variable QoS} & \textbf{Unidad de medida} & \textbf{Por qué es importante} \\
\hline
Throughput (Ancho de Banda Efectivo) & Gbps / Mbps & Determina cuánta información puede transmitirse por segundo. Herramienta: iperf3 -J, 300 s TCP. \\
Latencia de Establecimiento TCP & ms (min/avg/max/mdev) & Tiempo que tarda una conexión en establecerse. Herramienta: script Bash con \texttt{nc -z -w 2}, 30 muestras por corrida. \\
Retransmisiones TCP & Número de paquetes & Indicador de problemas en el canal de encapsulamiento del CNI. \\
Costo Computacional del CNI & \% CPU / MB RAM & Cuántos recursos del servidor consume el propio CNI para procesar el tráfico. \\
\hline
\end{tabularx}
\caption{Variables telemáticas capturadas por el sistema de medición y su impacto en las aplicaciones.}
\end{table*}

\subsubsection{La Pila de Observabilidad: Prometheus y Grafana}

\textbf{Prometheus:} Sistema de recolección de métricas. Funciona como un inspector que visita periódicamente cada nodo del clúster y registra estadísticas de CPU, memoria, red y el estado de cada CNI. Los datos se almacenan en series temporales, lo que permite correlacionar el consumo de recursos del agente CNI con los instantes exactos en que se ejecutan las pruebas de throughput y verificar que las métricas de eficiencia reflejen el comportamiento bajo carga y no en reposo \cite{ref6}.

\textbf{Grafana:} Es la capa de visualización. Conecta con Prometheus y presenta los datos en dashboards interactivos que muestran en tiempo real: el throughput de cada CNI, la latencia promedio, el consumo de CPU del agente CNI y las retransmisiones TCP. Esta visibilidad permitió detectar anomalías durante las sesiones de prueba ---como picos de CPU o reinicios inesperados del agente--- antes de que comprometieran la integridad de las corridas, y sirve como evidencia verificable del comportamiento del sistema \cite{ref6}.

\subsection{Diseño del Modelo Matemático de Recomendación (MCDA)}

Para resolver el problema de selección objetiva, se diseñó un modelo de Análisis de Decisión Multicriterio (MCDA). Este tipo de modelos se usan en ingeniería y economía cuando se deben tomar decisiones que involucran múltiples factores con diferentes niveles de importancia \cite{ref9}. A diferencia de una comparación basada en un único criterio ---como el throughput máximo---, el MCDA captura el trade-off inherente entre velocidad, eficiencia de recursos y capacidades de seguridad, produciendo una recomendación que refleja las prioridades reales de cada tipo de organización.

\subsubsection{Algoritmo de Procesamiento y Limpieza Estadística}

Antes de aplicar el modelo MCDA, los datos recolectados pasan por un proceso de limpieza estadística mediante el Rango Intercuartílico (IQR) para eliminar outliers causados por eventos externos (mantenimientos del proveedor, picos de uso compartido). Posteriormente, las métricas se normalizan con min-max a una escala común de 0 a 1, donde 1 representa el mejor desempeño posible. Este pipeline se implementa en el archivo \texttt{docs/procesador.js}.

\subsubsection{La Matriz de Decisión por Perfil de Usuario}

Se diseñaron tres perfiles de usuario basados en los patrones más comunes de uso de Kubernetes en la industria:

\begin{table*}[htbp]
\centering
\scriptsize
\begin{tabularx}{\textwidth}{|Y|Y|Y|Y|}
\hline
\textbf{Criterio de Evaluación} & \textbf{Fintech (Seguridad)} & \textbf{Streaming (Rendimiento)} & \textbf{IoT (Recursos)} \\
\hline
Throughput (velocidad) & 15\% & 35\% & 25\% \\
Latencia (tiempo de respuesta) & 20\% & 35\% & 20\% \\
Retransmisiones TCP (estabilidad) & 15\% & 20\% & 15\% \\
Consumo de CPU del CNI & 10\% & 10\% & 30\% \\
Consumo de RAM del CNI & 10\% & 0\% & 10\% \\
Soporte de Network Policies L7 & 30\% & 0\% & 0\% \\
Total & 100\% & 100\% & 100\% \\
\hline
\end{tabularx}
\caption{Matriz de pesos del modelo MCDA por perfil de usuario.}
\end{table*}

La ecuación general del modelo se expresa como:
\[
P_{CNI}=\sum_{i=1}^{n}(V_i \times W_i)
\]
donde $P_{CNI}$ es el puntaje total del plugin, $V_i$ es la valoración normalizada de la métrica $i$, $W_i$ es el peso asignado y $n$ es el número de métricas evaluadas.

\subsection{Diseño de Seguridad y Micro-segmentación (Modelo Zero Trust)}

En su configuración inicial, cualquier aplicación dentro del clúster puede comunicarse libremente con cualquier otra. Las Network Policies permiten implementar Zero Trust: permitir sólo lo necesario y verificar explícitamente cada solicitud de acceso \cite{ref24,ref25}. Esto reduce drásticamente la superficie de ataque: si un atacante compromete un Pod, las políticas de red limitan su capacidad de movimiento lateral a únicamente las rutas explícitamente autorizadas, convirtiendo el clúster en una red de confianza mínima en lugar de un entorno totalmente abierto.

\subsubsection{Estrategia de Default Deny}

El diseño implementa Zero Trust en dos niveles. El primer nivel es la política de Default Deny: se crean reglas que bloquean todo el tráfico entrante y saliente de cada namespace. El segundo nivel es la habilitación selectiva: una vez bloqueado todo, se abren únicamente las comunicaciones necesarias.

Se implementaron tres casos de uso ejecutados sobre cada CNI que soporta Network Policies (Calico, Cilium y Antrea); Flannel se incluye como referencia negativa para confirmar la ausencia de \emph{enforcement} en su plano de datos:

\begin{table*}[htbp]
\centering
\scriptsize
\begin{tabularx}{\textwidth}{|Y|Y|}
\hline
\textbf{Caso} & \textbf{Descripción} \\
\hline
\texttt{zero\_trust} & Default Deny: bloquea todo ingress y egress por defecto; permite solo el tráfico benchmark explícito. \\
\texttt{multi\_tier} & Reglas explícitas \texttt{frontend $\rightarrow$ backend $\rightarrow$ database} mediante label selectors. Cilium agrega CiliumNetworkPolicy L7 HTTP (solo GET/HEAD desde frontend). \\
\texttt{egress\_block} & Bloqueo de egress externo con DNS permitido; tráfico inter-pod interno habilitado. \\
\hline
\end{tabularx}
\caption{Casos de uso de seguridad implementados por CNI.}
\end{table*}

\textbf{Diferenciadores únicos por CNI:}
\begin{itemize}
  \item \textbf{Cilium} (\texttt{CiliumNetworkPolicy} L7): filtrado HTTP nativo por método (\texttt{GET}/\texttt{HEAD} permitido desde frontend a backend; \texttt{POST}/\texttt{PUT}/\texttt{DELETE} bloqueados) directamente en eBPF, sin service mesh.
  \item \textbf{Antrea} (\texttt{ClusterNetworkPolicy} en Tier \texttt{securityops}): las reglas en el tier SecurityOps tienen prioridad sobre cualquier \texttt{NetworkPolicy} de namespace, modelando políticas corporativas obligatorias (ISO 27001 / SOC2) que los desarrolladores no pueden sobrescribir.
\end{itemize}

\subsubsection{Automatización GitOps con ArgoCD}

Para garantizar que las políticas de red sigan siendo efectivas después de cambios en el sistema, se diseñó un flujo de trabajo GitOps usando ArgoCD. Si alguien elimina una política de red del clúster, ArgoCD detecta la diferencia entre el estado actual y el estado definido en Git, y automáticamente restaura la política eliminada (\texttt{selfHeal: true}). Los recursos eliminados del repositorio también se eliminan del clúster (\texttt{prune: true}). Esta automatización garantiza además la integridad del experimento: durante las rotaciones de CNI, las políticas de red se restauran a su estado esperado sin intervención manual, eliminando una fuente potencial de sesgo en los resultados comparativos \cite{ref25}.

\subsection{Flujo Integral del Sistema}

El flujo completo de los datos tiene cinco etapas:
\begin{enumerate}
\item \textbf{Generación del tráfico:} CronJob activa iperf3-client cada 30 minutos (minutos 0 y 30). La anti-afinidad ubica cliente y servidor en nodos distintos.
\item \textbf{Captura de métricas:} iperf3 registra throughput, latencia TCP y retransmisiones (formato JSON, 300 s). CronJob de latencia ejecuta 30 muestras con \texttt{nc -z -w 2} (minutos 10 y 40). Prometheus recolecta CPU y RAM del agente CNI.
\item \textbf{Extracción de recursos:} CronJob del Resource Usage Exporter (minutos 25 y 55) consulta directamente la API HTTP de Prometheus (\texttt{/api/v1/query\_range}) y genera archivos raw JSON, CSV, PNG (gnuplot) y \texttt{cni\_resource\_summary.json}.
\item \textbf{Limpieza, normalización y scoring MCDA:} \texttt{procesador.js} aplica IQR, normalización min-max y calcula puntajes. Genera \texttt{cni-data.json}.
\item \textbf{Visualización y recomendación:} SPA React consume \texttt{cni-data.json} y muestra el ranking MCDA por perfil.
\end{enumerate}

\subsection{Beneficio Tecnológico y Telemático de la Solución}

\subsubsection{Beneficios Telemáticos}
El sistema proporciona visibilidad completa de la red en tiempo real y control de seguridad mediante micro-segmentación Zero Trust, transformando la red de un sistema ``todo abierto'' a una red controlada. Esto reduce la superficie de ataque: si un atacante compromete una aplicación, su capacidad de movimiento lateral está limitada por las políticas de red \cite{ref16}.

\subsubsection{Beneficios Tecnológicos}
Reemplaza las decisiones basadas en opiniones por decisiones basadas en datos medidos en condiciones controladas y reproducibles. El modelo MCDA proporciona un método objetivo y transparente, y la metodología es replicable para evaluar nuevas versiones de CNIs.

\subsubsection{Beneficio Económico: Reducción del Sobredimensionamiento}
Con las mediciones del proyecto, es posible saber exactamente cuánta CPU consume cada CNI para mover 1 Gbps de tráfico. Esta información permite un dimensionamiento inteligente: si Cilium necesita un 30\% menos de CPU que Flannel para la misma cantidad de tráfico, una organización podría reducir su inversión en cómputo de manera proporcional.

\subsection{Justificación de las Principales Decisiones de Diseño}

\begin{table*}[htbp]
\centering
\scriptsize
\begin{tabularx}{\textwidth}{|Y|Y|}
\hline
\textbf{Decisión de diseño} & \textbf{Justificación técnica y telemática} \\
\hline
¿Por qué nube pública (DigitalOcean) en lugar de un laboratorio físico propio? & Los entornos físicos propios introducen variables no controlables. La nube garantiza que la única variable entre pruebas sea el CNI, no el hardware. \\
¿Por qué K3s y no Kubernetes estándar? & K3s tiene un footprint menor, lo que hace que una mayor proporción del CPU y la RAM esté disponible para las aplicaciones y el CNI, aumentando la sensibilidad de las mediciones. \\
¿Por qué iperf3? & Es la herramienta de referencia académica y profesional para pruebas de rendimiento de redes TCP/UDP. Sus resultados son reproducibles y comparables con otros estudios. \\
¿Por qué IQR para limpiar outliers? & El promedio es muy sensible a valores extremos. El IQR es robusto a outliers porque se basa en la distribución estadística central. \\
¿Por qué tres perfiles de usuario en el MCDA? & Existe una tensión fundamental entre throughput, latencia, seguridad y uso de recursos. Ningún CNI es el mejor en todas las dimensiones simultáneamente. \\
¿Por qué ArgoCD para la gestión de políticas? & Las políticas de red manuales son frágiles. GitOps con ArgoCD garantiza que el estado de seguridad sea siempre el que el equipo definió, sin importar qué cambios ocurran en el clúster \cite{ref25}. \\
\hline
\end{tabularx}
\caption{Justificación de las principales decisiones de diseño del proyecto.}
\end{table*}

\subsection{Resumen Arquitectónico Integrado}

El diseño de la solución puede resumirse como un sistema de tres capas:

La \textbf{primera capa} es la infraestructura controlada (red underlay sobre DigitalOcean con Terraform), que garantiza un entorno reproducible y libre de ruido externo.

La \textbf{segunda capa} es el plano de observabilidad y prueba (red overlay con los cuatro CNIs, iperf3 con podAntiAffinity, CronJobs automáticos, Prometheus y Grafana), que captura datos de calidad de servicio de manera sistemática.

La \textbf{tercera capa} es la inteligencia de decisión (algoritmo IQR, normalización, modelo MCDA con perfiles, dashboard de recomendación), que transforma los datos brutos en recomendaciones específicas.

Transversalmente opera la estrategia de seguridad Zero Trust (Network Policies con Default Deny, micro-segmentación y automatización GitOps con ArgoCD), que garantiza que el entorno de prueba y las recomendaciones incluyan siempre consideraciones de seguridad.

\subsection{Trazabilidad del diseño frente al problema de investigación}

El diseño de la solución no se planteó como una arquitectura aislada, sino como una respuesta directa al problema formulado en el anteproyecto: las organizaciones adoptan Kubernetes como plataforma de orquestación, pero la calidad final del despliegue depende en gran medida del plano de red, de la selección del CNI y de la forma en que se controlan las comunicaciones internas del clúster \cite{ref29}.

\begin{table*}[htbp]
\centering
\scriptsize
\caption{Trazabilidad entre problema de investigación, objetivo y decisión de diseño.}
\label{tab:trazabilidad-diseno}
\begin{tabularx}{\textwidth}{|Y|Y|Y|}
\hline
\textbf{Necesidad identificada} & \textbf{Objetivo asociado} & \textbf{Respuesta de diseño} \\
\hline
Comparar CNIs sin depender de recomendaciones informales. & Evaluar cuantitativamente Calico, Cilium, Flannel y Antrea en escenarios controlados. & Testbed reproducible sobre Kubernetes, ejecución automatizada y consolidación de métricas por CNI. \\
\hline
Controlar el riesgo de comunicaciones internas abiertas. & Evaluar políticas de red seguras y automatizadas. & Diseño de NetworkPolicies bajo enfoque Zero Trust, casos de denegación por defecto y segmentación por capas. \\
\hline
Convertir resultados técnicos heterogéneos en una decisión comprensible. & Definir criterios, umbrales y ponderaciones para comparación objetiva. & Modelo MCDA con criterios de rendimiento, eficiencia de recursos e impacto de seguridad. \\
\hline
Apoyar a una organización que necesita seleccionar y configurar un CNI. & Implementar un prototipo funcional de recomendación. & Prototipo web que consume datos procesados, aplica perfiles de decisión y explica la recomendación. \\
\hline
Evitar sobredimensionamiento de infraestructura. & Relacionar desempeño de red con consumo de CPU y memoria. & Captura de recursos del CNI mediante Prometheus y cálculo de eficiencia relativa por perfil. \\
\hline
\end{tabularx}
\end{table*}

\subsection{Requisitos de diseño derivados del contexto telemático}

El primer requisito es la \textbf{reproducibilidad}. El entorno se concibió bajo infraestructura como código, despliegue declarativo y automatización de experimentos. El segundo requisito es el \textbf{aislamiento de variables}: el CNI es la única variable; todos los demás elementos (proveedor cloud, tipo de instancia, versión de Kubernetes, topología, duración) se mantienen constantes. El tercer requisito es la \textbf{medición integral}: un CNI no debe evaluarse solo por throughput, sino por rendimiento de red, eficiencia de recursos e impacto de seguridad. El cuarto requisito es la \textbf{trazabilidad de la evidencia}: cada dato debe conservar su origen. El quinto requisito es la \textbf{compatibilidad con la operación real}: el diseño usa objetos y patrones nativos de Kubernetes (Deployments, CronJobs, ArgoCD) cercanos a la forma en que los equipos DevOps y SRE operan clústeres reales.

\subsection{Diseño arquitectónico por capas}

La solución se organizó en siete capas para separar responsabilidades. La \textbf{capa 1} es la infraestructura underlay (DigitalOcean). La \textbf{capa 2} es el clúster Kubernetes K3s en alta disponibilidad. La \textbf{capa 3} es el plano CNI (Flannel, Calico, Cilium o Antrea). La \textbf{capa 4} es el plano de inyección de tráfico (Pods iperf3 + CronJobs). La \textbf{capa 5} es observabilidad y captura (Prometheus, Grafana, exporter de recursos). La \textbf{capa 6} es procesamiento y decisión (\texttt{procesador.js}: IQR, normalización, MCDA). La \textbf{capa 7} es presentación y recomendación (SPA React).

\subsection{Diseño comparativo de los CNIs evaluados}

La selección de Flannel, Calico, Cilium y Antrea no es accidental. Estos cuatro CNIs representan enfoques distintos del plano de datos. Flannel representa simplicidad y conectividad básica (línea base). Calico representa políticas de red y ruteo con fuerte adopción en producción. Cilium representa el uso de eBPF como mecanismo moderno de aceleración y observabilidad. Antrea representa una aproximación basada en Open vSwitch y conceptos de redes definidas por software.

\begin{table*}[htbp]
\centering
\scriptsize
\caption{Criterios de diseño para la selección de CNIs evaluados.}
\label{tab:cni-diseno-comparativo}
\begin{tabularx}{\textwidth}{|Y|Y|Y|Y|}
\hline
\textbf{CNI} & \textbf{Enfoque técnico} & \textbf{Aporte al experimento} & \textbf{Riesgo o limitación a medir} \\
\hline
Flannel & Overlay simple (VXLAN). & Línea base de conectividad sencilla. & Ausencia de enforcement nativo de NetworkPolicies. \\
\hline
Calico & Ruteo y políticas de red; soporte de controles avanzados. & Comparación de seguridad práctica y madurez operativa. & Posible costo por evaluación de reglas. \\
\hline
Cilium & Dataplane y seguridad apoyados en eBPF. & Evaluar beneficios de filtrado moderno, visibilidad y políticas L7. & Mayor complejidad conceptual y alto consumo de RAM. \\
\hline
Antrea & Dataplane basado en Open vSwitch y enfoque SDN. & Evaluar conmutación virtual programable y visibilidad de flujos. & Dependencia de componentes adicionales. \\
\hline
\end{tabularx}
\end{table*}

\subsection{Diseño de la metrología de QoS}

El sistema de medición se diseñó bajo el principio de que una métrica aislada no describe la calidad de una red. El throughput TCP mide la capacidad efectiva de transferencia. La latencia de establecimiento TCP representa el costo de iniciar una comunicación entre servicios. El jitter mide la estabilidad temporal. Las retransmisiones TCP son un indicador de eficiencia del dataplane. El consumo de CPU y memoria determina el impacto en el sobredimensionamiento.

\begin{table*}[htbp]
\centering
\scriptsize
\caption{Métricas de QoS y eficiencia incluidas en el diseño.}
\label{tab:metricas-diseno-qos}
\begin{tabularx}{\textwidth}{|Y|Y|Y|}
\hline
\textbf{Métrica} & \textbf{Qué observa} & \textbf{Uso dentro del modelo} \\
\hline
Throughput TCP & Capacidad efectiva de transferencia entre Pods. & Determinar eficiencia del dataplane para cargas de alto volumen. \\
\hline
Latencia TCP connect & Tiempo de establecimiento de conexión. & Evaluar respuesta para microservicios y tráfico transaccional. \\
\hline
Latencia máxima y variabilidad & Estabilidad temporal de la comunicación. & Identificar picos que el promedio puede ocultar. \\
\hline
Retransmisiones TCP & Señal de reintentos, congestión o inestabilidad del camino. & Penalizar dataplanes que logren throughput a costa de más retransmisiones. \\
\hline
CPU del agente CNI & Costo computacional del plano de red. & Medir impacto en sobredimensionamiento y perfiles edge. \\
\hline
Memoria del agente CNI & Huella de memoria persistente por nodo. & Evaluar sostenibilidad operativa en nodos restringidos. \\
\hline
Overhead con NetworkPolicies & Degradación al activar seguridad. & Comparar seguridad aplicable sin comprometer rendimiento. \\
\hline
\end{tabularx}
\end{table*}

\subsection{Diseño estadístico y modelo de decisión multicriterio}

El procesamiento estadístico se diseñó para evitar que la recomendación dependa de valores atípicos. Después de la limpieza IQR, las métricas se normalizan considerando el sentido de optimización de cada variable: en throughput, un valor mayor es mejor; en latencia, CPU, memoria y retransmisiones, un valor menor es preferible (la escala se invierte). La decisión final se formaliza con la suma ponderada del MCDA.

El diseño contempla tres perfiles arquitectónicos: URLLC/automatización industrial (prioriza latencia y jitter), Edge/IoT (prioriza consumo de recursos y footprint) y microservicios transaccionales/Zero Trust (prioriza seguridad y bajo overhead con políticas activas). Los pesos son parámetros explícitos del modelo, lo que permite que el prototipo sea ajustable por organización.

\subsection{Diseño de seguridad: micro-segmentación y Zero Trust}

La seguridad del diseño parte de una característica esencial de Kubernetes: la comunicación interna entre Pods suele estar abierta por defecto. El diseño adopta Zero Trust como principio operativo \cite{ref24,ref25}: no se asume confianza por pertenecer al mismo clúster o namespace. Las NetworkPolicies dependen del CNI para su enforcement real; si el plugin no implementa el plano de datos correctamente, la política puede existir como objeto declarativo sin producir el efecto esperado \cite{ref4,ref18}.

La comparación de seguridad no se limita a verificar si una política funciona. También evalúa capacidades diferenciadoras: Flannel como referencia sin soporte nativo, Calico con políticas estándar y extensiones, Cilium con capacidades L7 mediante CiliumNetworkPolicy, y Antrea con estructuras de prioridad y visibilidad asociadas a OVS.

\subsection{Diseño GitOps para consistencia operativa}

La automatización GitOps cumple dos funciones. La primera es operativa: garantizar que los manifiestos aplicados al clúster correspondan al estado versionado en el repositorio. La segunda es metodológica: permitir que las pruebas sean repetibles y auditables. Un cambio de CNI o de escenario de prueba se expresa como un cambio declarativo; el clúster sincroniza ese estado, ejecuta las cargas programadas y produce resultados en la estructura de datos esperada.

\subsection{Diseño del prototipo recomendador}

El prototipo recomendador se diseñó como la capa de interacción entre el modelo matemático y el usuario final. La salida esperada incluye tres niveles: (1) la recomendación directa (qué CNI se sugiere), (2) la justificación (qué métricas explican la recomendación) y (3) la orientación de configuración (qué consideraciones debe tener el usuario al desplegar ese CNI). Esta estructura responde al objetivo específico 4 del anteproyecto \cite{ref29}.

El prototipo se implementó como SPA con React~19, Vite~7 y Tailwind CSS~3. El build de producción se despliega en GitHub Pages vía \texttt{build:pages} (\texttt{--outDir ../docs/recommender}). Para actualizar datos: \texttt{npm run sync:data} ejecuta \texttt{procesador.js} (genera \texttt{cni-data.json}) y luego \texttt{export-spa-data.cjs} (inyecta los datos en el bundle).

\subsection{Alcance, límites y decisiones explícitas del diseño}

El diseño se limita a evaluar CNIs en un testbed controlado sobre Kubernetes y no pretende generalizar automáticamente los resultados a cualquier proveedor, versión o hardware. Se usa tráfico TCP controlado como señal principal. El diseño tampoco asume que mayor complejidad implica mejor resultado. La solución se diseñó para sostener una mejora continua: a medida que se agreguen nuevas corridas o nuevas versiones de CNIs, el pipeline puede recalcular puntajes. En lugar de entregar una decisión congelada, entrega una forma de decidir con evidencia actualizable \cite{ref29}.

\section{Implementación de la Solución}

Este capítulo documenta cómo se ejecutó el diseño del capítulo anterior. Toda la infraestructura está codificada en Terraform, toda la configuración en manifiestos YAML versionados en Git, y toda prueba está automatizada en CronJobs de Kubernetes ---condición necesaria para que los resultados sean reproducibles y atribuibles al CNI \cite{ref7}.

\subsection{Aprovisionamiento de Infraestructura con Terraform}

El aprovisionamiento de servidores y redes se realizó de forma completamente automatizada mediante Terraform (IaC). Al describir la infraestructura en archivos versionados, Terraform la recrea de manera idéntica en cada ejecución y elimina la posibilidad de que variaciones manuales en el aprovisionamiento introduzcan diferencias entre los entornos de cada CNI \cite{ref7}.

\subsubsection{Estructura del Código Terraform (K8s-bootstrap/00.bootstrap)}

% [EDIT] Tabla booktabs y caption superior.
\begin{table*}[htbp]
\centering
\scriptsize
\renewcommand{\arraystretch}{1.25}
\setlength{\tabcolsep}{4pt}
\caption{Archivos Terraform del módulo de aprovisionamiento y su función específica.}
\label{tab:terraform-aprovisionamiento}
\begin{tabularx}{\textwidth}{|l|Y|}
\hline
\textbf{Archivo Terraform} & \textbf{Responsabilidad en la Infraestructura} \\
\hline
\texttt{main.tf} & Define los recursos principales: Droplets de DigitalOcean y red VPC privada. \\
\hline
\texttt{variables.tf} & Centraliza todos los parámetros configurables: número de nodos (\texttt{agent\_count = 2}), tamaño de instancias, región y proveedor CNI. \\
\hline
\texttt{ssh\_keys.tf} & Automatiza la gestión de llaves SSH para el acceso seguro entre nodos. \\
\hline
\texttt{firewall.tf} & Configura las reglas de firewall a nivel de nube: primera línea de seguridad perimetral. \\
\hline
\texttt{outputs.tf} & Genera automáticamente el archivo kubeconfig con las credenciales del clúster K3s. \\
\hline
\end{tabularx}
\end{table*}

El proceso de aprovisionamiento tiene tres etapas: \texttt{terraform plan} (compara estado actual vs. código), \texttt{terraform apply} (crea Droplets, VPC y firewall) y \texttt{terraform output} (extrae IPs y kubeconfig). Terraform mantiene un archivo \texttt{terraform.tfstate} que permite destruir y recrear toda la infraestructura de manera controlada, garantizando la repetibilidad científica: cada CNI se evaluó en condiciones idénticas de arranque.

\subsection{Configuración del Clúster K3s en Alta Disponibilidad}

Una vez aprovisionados los servidores, se instaló y configuró K3s, una distribución liviana de Kubernetes, con topología de 1 nodo Master + 2 nodos Worker dedicados.

\subsubsection{Base de Datos Externa para el Plano de Control}

El plano de control se configuró con una base de datos PostgreSQL gestionada de DigitalOcean (también provisionada via Terraform) como datastore externo. Esto elimina la complejidad de mantener un clúster etcd y garantiza un plano de control sin punto único de fallo (SPOF).

\subsubsection{Instalación: clúster sin CNI para plugins externos y modo nativo para Flannel}

La implementación establece una distinción metodológica deliberada: Flannel se evalúa en su modo de despliegue más habitual ---integrado nativamente en K3s con \texttt{flannel-iface: eth1}---, mientras que Calico, Cilium y Antrea se instalan como plugins externos sobre un clúster sin CNI. La asimetría refleja las condiciones reales de uso: Flannel es el CNI predeterminado de K3s y la mayoría de equipos lo adopta sin configuración adicional, en tanto que los otros tres requieren instalación explícita y declarativa. Evaluar cada CNI en su escenario de despliegue típico aumenta la validez ecológica del experimento; debe considerarse, no obstante, que Flannel se beneficia de la integración nativa de K3s ---gestión automática de CIDR y \emph{hooks} de kubelet--- que los plugins externos deben replicar por su cuenta.

Para los CNIs externos, la configuración crítica es el parámetro \texttt{--flannel-backend=none}, que desactiva completamente el CNI integrado de K3s y evita interferencias con el plugin bajo evaluación. Adicionalmente, se configura \texttt{--disable-network-policy} para que K3s no intente implementar NetworkPolicies por sí mismo.

La configuración resultante en \texttt{/etc/rancher/k3s/config.yaml} para CNIs externos incluye:
\begin{lstlisting}
flannel-backend: none
disable-network-policy: true
\end{lstlisting}

% [EDIT] Tabla booktabs y label consistente.
\begin{table*}[htbp]
\centering
\scriptsize
\renewcommand{\arraystretch}{1.25}
\setlength{\tabcolsep}{4pt}
\caption{Instalación de CNIs: método y configuración clave en \texttt{user\_data/ks3\_server\_init.sh}.}
\label{tab:cni-instalacion}
\begin{tabularx}{\textwidth}{|l|p{3.2cm}|Y|}
\hline
\textbf{CNI} & \textbf{Método de instalación} & \textbf{Configuración clave} \\
\hline
Flannel & Nativo K3s & \texttt{flannel-iface: eth1} (interfaz VPC de DigitalOcean) \\
\hline
Calico & Tigera Operator v3.29.1 & \texttt{encapsulation: VXLAN}, CIDR \texttt{10.42.0.0/16}, \texttt{nodeAddressAutodetectionV4.interface: eth1} \\
\hline
Cilium & Helm 1.16.5 & \texttt{routingMode=tunnel}, \texttt{tunnelProtocol=vxlan}, \texttt{ipam.mode=kubernetes}, \texttt{devices=eth1}, \texttt{operator.replicas=1} \\
\hline
Antrea & Manifest v2.2.0 & \texttt{transportInterface: eth1}, \texttt{serviceCIDR: 10.43.0.0/16} \\
\hline
\end{tabularx}
\end{table*}

El protocolo de rotación de CNIs tuvo cuatro pasos: (1) \texttt{kubectl delete} para eliminar recursos del CNI actual; (2) verificación de que no quedan interfaces de red anteriores en los nodos; (3) reinicio de los nodos Worker para limpiar estado en memoria; (4) aplicación del manifiesto del nuevo CNI via ArgoCD. Este protocolo garantiza que cada CNI inicia sobre un nodo completamente limpio, sin residuos de configuraciones anteriores que pudieran contaminar las métricas del siguiente ciclo de pruebas.

\subsection{Orquestación y Entrega Continua: GitOps con ArgoCD}

Con el clúster K3s funcionando y sin CNI, el siguiente componente instalado fue ArgoCD. ArgoCD implementa el paradigma GitOps: el repositorio Git es la única fuente de verdad sobre el estado del clúster. Todo lo que debe existir en el clúster está definido en archivos YAML dentro del repositorio \texttt{K8s-Tesis-Apps}. Esta centralización en Git facilita la auditoría del experimento: cualquier cambio en la configuración queda registrado con autor, fecha y motivo, haciendo el proceso de prueba completamente trazable.

\subsubsection{El Patrón App of Apps}

La implementación de ArgoCD utiliza el patrón App of Apps: existe UNA sola aplicación raíz (ubicada en \texttt{K8s-bootstrap/30.apps-of-apps-install}) que ArgoCD conoce directamente, y esa aplicación raíz contiene las definiciones de todas las demás aplicaciones (CNIs, stack de métricas, benchmarks, políticas de red).

\subsubsection{Configuración de selfHeal y prune}

\textbf{\emph{selfHeal: true}} --- Ante la modificación manual de un recurso del clúster ---por ejemplo, la eliminación de una NetworkPolicy---, ArgoCD detecta la diferencia con el repositorio Git y restaura el estado original, lo que resultó crítico para la integridad del experimento.

\textbf{\emph{prune: true}} --- Al eliminar un archivo YAML del repositorio Git, ArgoCD elimina también el recurso correspondiente del clúster, lo que evita la acumulación de recursos huérfanos.

\begin{lstlisting}
# Definicion de la Application raiz en ArgoCD (App of Apps)
apiVersion: argoproj.io/v1alpha1
kind: Application
metadata:
  name: apps-of-apps
  namespace: argocd
spec:
  source:
    repoURL: https://github.com/equipo/K8s-Tesis-Apps
    targetRevision: main
    path: apps/
  syncPolicy:
    automated:
      selfHeal: true   # Restaura cambios manuales automaticamente
      prune: true      # Elimina recursos borrados del repo
\end{lstlisting}

\subsubsection{Overlays de Kustomize para los CNIs}

Cada CNI tiene una configuración ligeramente diferente. Para gestionarlas sin duplicar código, se utilizó Kustomize con el patrón base/overlay: existe una configuración base con los parámetros comunes, y sobre esa base se aplican overlays específicas para cada CNI.

Las apps gestionadas por ArgoCD incluyen: \texttt{cni-test-iperf} (benchmarks de throughput y latencia), \texttt{cni-resource-exporter} (exportación de recursos desde Prometheus), \texttt{cni-np-test} (pruebas de Network Policies, creada solo cuando \texttt{cni\_provider != "flannel"} y el caso de uso está definido), y \texttt{metrics} (Prometheus + Grafana).

\subsection{Implementación del Stack de Observabilidad}

El stack de observabilidad se implementó mediante el \texttt{kube-prometheus-stack} (Helm chart) que instala Prometheus, Grafana y exportadores preconfigurados para Kubernetes.

\subsubsection{Prometheus: El Recolector de Métricas}

Prometheus opera con un modelo pull: visita periódicamente los endpoints de métricas (cada 15 segundos en la configuración del proyecto). Para que Prometheus sepa qué métricas recolectar de los agentes CNI, se implementaron \texttt{ServiceMonitors} (CRD), apuntando específicamente al endpoint de métricas de cada agente. Un intervalo de 15 segundos proporciona suficiente granularidad para detectar picos de consumo durante las ráfagas de tráfico iperf3, sin generar una carga de scraping que interfiera con las métricas bajo observación.

Las métricas clave recolectadas del agente CNI incluyen:
\begin{itemize}
\item \path|container_cpu_usage_seconds_total| (filtrado por pod del agente): costo computacional del CNI (criterio C2 del marco MCDA).
\item \path|container_memory_rss| (filtrado por pod del agente): memoria física real del proceso del agente CNI.
\item \path|node_network_transmit_bytes_total| / \path|node_network_receive_bytes_total|: bytes totales de red por nodo.
\item \path|kube_pod_info| + \path|kube_pod_status_phase|: estado del agente CNI para detectar reinicios durante las pruebas.
\end{itemize}

\subsubsection{Grafana: Los Dashboards de Red}

% latex-paper-en · deai (Line 142) [Severity: P0] [Priority: P0]: lista paralela mecánica
Grafana centraliza la lectura operativa de las series de Prometheus. Sus paneles agrupan throughput, latencia, consumo del agente CNI y score MCDA. Esta correlación vincula ráfagas iperf3 con CPU, RSS y el ranking calculado por \texttt{procesador.js}.

\subsubsection{Resource Usage Exporter: Extracción Automatizada de Datos}

% latex-paper-en · deai (Line 147) [Severity: P1] [Priority: P1]: tono operativo repetitivo
El \texttt{cni-resource-usage-exporter} corre como CronJob en \path|K8s-Tesis-Apps/grafana-exporter|. ArgoCD sincroniza el recurso mediante la app \texttt{cni-resource-exporter}. El CronJob:
\begin{itemize}
\item Ejecuta la extracción en los minutos 25 y 55 de cada hora, tras cada ventana de benchmarks.
\item Consulta directamente la API HTTP de Prometheus (\texttt{/api/v1/query\_range}) --- no interactúa con Grafana.
\item Recolecta CPU\% por nodo, RAM\% por nodo, RAM usada, CPU cores y RAM del agente CNI.
\item Genera \texttt{raw/}, \texttt{csv/}, \texttt{images/} y \texttt{cni\_resource\_summary.json} por corrida.
\item Persiste cada salida con timestamp en \path|K8S-CNI-Results/results/cni-benchmarks/{cni}/resource_usage_nodes/|.
\end{itemize}

\subsection{Implementación del Banco de Pruebas de QoS}

\subsubsection{Escenario de Rendimiento con iperf3}

% latex-paper-en · deai (Line 161) [Severity: P1] [Priority: P1]: frase genérica
El banco de throughput se define en \texttt{K8s-Tesis-Apps/cni\_test\_iperf}:

\textbf{El iperf3-server (Deployment):} Un Deployment ejecuta una réplica de iperf3 en modo servidor. Un Service ClusterIP expone un nombre DNS estable para el cliente.

\textbf{El iperf3-client (CronJob):} La programación \texttt{0,30 * * * *} activa el cliente en los minutos 0 y 30. \texttt{concurrencyPolicy: Forbid} bloquea corridas simultáneas. Cada activación crea un Pod y genera tráfico TCP durante 300 segundos. La salida conserva formato JSON con \texttt{iperf3 -t 300 -i 10 -J}. La retención máxima queda en 5 corridas (\texttt{RETENTION\_MAX\_RUNS="5"}).

La anti-afinidad del cliente frente al servidor se implementó como \texttt{preferredDuringSchedulingIgnoredDuringExecution} (suave, weight 100), forzando tráfico inter-nodo en la práctica. Esta ubicación garantiza que todo el tráfico de benchmark atraviese el mecanismo de encapsulación del CNI entre nodos, que es exactamente lo que se quiere medir.

\begin{lstlisting}
# iperf-client-cronjob.yaml -- fragmento relevante
schedule: '0,30 * * * *'   # Minutos 0 y 30
concurrencyPolicy: Forbid
affinity:
  podAntiAffinity:
    preferredDuringSchedulingIgnoredDuringExecution:
    - weight: 100
      podAffinityTerm:
        labelSelector:
          matchLabels:
            app: iperf-server
        topologyKey: kubernetes.io/hostname
# iperf3 -c iperf-server-svc -t 300 -i 10 -J
\end{lstlisting}

\subsubsection{Cliente de Latencia Personalizado}

% latex-paper-en · deai (Line 188) [Severity: P1] [Priority: P1]: frase temporal mecánica
El CronJob \texttt{latency-client-cronjob.yaml} ejecuta la prueba de latencia. Su programación corre en los minutos 10 y 40 de cada hora. Esa separación evita solapamiento con el inicio del throughput.

El cliente registra TCP Connect con \texttt{nc -z -w 2} y timestamp en ms. Cada corrida toma 30 muestras y calcula min/avg/max/mdev. El JSON normalizado conserva campos de latencia, muestras intentadas y muestras exitosas.

% [EDIT] Tabla booktabs y label consistente.
\begin{table*}[htbp]
\centering
\scriptsize
\renewcommand{\arraystretch}{1.25}
\setlength{\tabcolsep}{4pt}
\caption{Métricas de QoS capturadas por el banco de pruebas y su relevancia telemática.}
\label{tab:metricas-qos-implementacion}
\begin{tabularx}{\textwidth}{|p{4cm}|p{2.4cm}|Y|}
\hline
\textbf{Métrica Capturada} & \textbf{Unidad} & \textbf{Relevancia para QoS Telemático} \\
\hline
Throughput TCP sostenido & Gbps / Mbps & Capacidad efectiva de la red. \\
\hline
Latencia mínima TCP connect & ms & Overhead mínimo del CNI sin congestión. \\
\hline
Latencia promedio TCP connect & ms & Comportamiento típico de solicitudes. \\
\hline
Latencia máxima / Jitter & ms & Peor caso y variabilidad. Crítico para tiempo real. \\
\hline
Retransmisiones TCP & count / sesión & Indicador de problemas en encapsulamiento o congestión. \\
\hline
CPU del agente CNI & millicores (m) & Impacto directo en costo de infraestructura cloud. \\
\hline
RAM del agente CNI (RSS) & MiB & Huella de memoria del agente CNI por nodo. \\
\hline
\end{tabularx}
\end{table*}

\subsection{Procesamiento Estadístico y Aplicación del Modelo MCDA}

\subsubsection{Limpieza con el Rango Intercuartílico (IQR)}

% latex-paper-en · deai (Line 226) [Severity: P2] [Priority: P2]: secuencia formulaica
\texttt{docs/procesador.js} transforma datos crudos en datos consolidados. El script descubre los CNIs disponibles bajo \texttt{results/cni-benchmarks}. Luego extrae métricas desde JSON y CSV. Finalmente aplica IQR:

\begin{lstlisting}
// procesador.js -- Funcion de limpieza estadistica IQR
function calcularEstadisticasLimpias(valores) {
  const sorted = [...valores].sort((a, b) => a - b);
  const q1 = sorted[Math.floor(sorted.length * 0.25)];
  const q3 = sorted[Math.floor(sorted.length * 0.75)];
  const iqr = q3 - q1;
  const lowerBound = q1 - 1.5 * iqr;
  const upperBound = q3 + 1.5 * iqr;
  const valoresLimpios = valores.filter(
    v => v >= lowerBound && v <= upperBound);
  const promedio = valoresLimpios.reduce(
    (sum, v) => sum + v, 0) / valoresLimpios.length;
  return { promedio,
    outliersPorcentaje: ((valores.length - valoresLimpios.length)
      / valores.length * 100).toFixed(1)
  };
}
\end{lstlisting}

% latex-paper-en · deai (Line 249) [Severity: P1] [Priority: P1]: causalidad débil
El pipeline descarta valores atípicos antes de calcular promedios. Dado que un evento cloud anómalo puede sesgar una corrida, el filtro estabiliza el puntaje MCDA del CNI.

\subsubsection{Normalización y Cálculo del Score MCDA}

% latex-paper-en · deai (Line 254) [Severity: P1] [Priority: P1]: estructura repetida
Después de la limpieza, el pipeline aplica normalización min-max. Cada valor se escala al rango $[0, 1]$. En métricas positivas, como throughput, la normalización es directa. En métricas negativas, como latencia o CPU, la escala se invierte.

Cuando la métrica penaliza, el valor más bajo recibe 1 y el más alto recibe 0. La escala común compara Gbps, milisegundos y millicores sin sesgo por magnitud numérica.

\begin{lstlisting}
// procesador.js -- Calculo del score MCDA por perfil
const perfiles = {
  fintech:   { throughput: 0.15, latencia: 0.20,
               retransmisiones: 0.15, cpu: 0.10,
               ram: 0.10, networkPoliciesL7: 0.30 },
  streaming: { throughput: 0.35, latencia: 0.35,
               retransmisiones: 0.20, cpu: 0.10,
               ram: 0.00, networkPoliciesL7: 0.00 },
  iot:       { throughput: 0.25, latencia: 0.20,
               retransmisiones: 0.15, cpu: 0.30,
               ram: 0.10, networkPoliciesL7: 0.00 }
};
\end{lstlisting}

% latex-paper-en · deai (Line 274) [Severity: P2] [Priority: P2]: transición sin carga lógica
El pipeline calcula overhead de seguridad cuando detecta \texttt{with\_network\_policy}. En ese caso, compara la métrica base contra la métrica con política activa.

\subsection{Prototipo Recomendador: Aplicación Web SPA}

\subsubsection{Tecnologías de Implementación}

La SPA fue implementada con React~19 (gestión de estado por componentes), Vite~7 (bundler; build de producción vía \texttt{build:pages} con \texttt{--outDir ../docs/recommender}) y Tailwind CSS~3 (estilos utilitarios directamente en JSX, sin CSS personalizado).

% latex-paper-en · deai (Line 305) [Severity: P2] [Priority: P2]: secuencia plana
\texttt{npm run sync:data} ejecuta \texttt{procesador.js} y genera \texttt{cni-data.json}. Después, \texttt{export-spa-data.cjs} inyecta los datos en el bundle. \texttt{recommendationModel.js} concentra la lógica de recomendación.

\subsubsection{Flujo de la Aplicación}

\begin{enumerate}
\item \textbf{Selección de Perfil:} La pantalla principal muestra tres tarjetas: Fintech/Seguridad, Streaming/Rendimiento e IoT/Recursos.
\item \textbf{Visualización de Resultados:} La aplicación carga \texttt{cni-data.json} y calcula el ranking MCDA para el perfil seleccionado. El CNI con mayor puntaje queda marcado como recomendación.
\item \textbf{Justificación Técnica:} La interfaz expone las métricas brutas promedio de cada CNI en gráficas comparativas.
\end{enumerate}

\subsection{Implementación de Network Policies y Seguridad Zero Trust}

\subsubsection{Políticas Default Deny}

% latex-paper-en · deai (Line 320) [Severity: P0] [Priority: P0]: paralelismo mecánico
Default Deny operacionaliza el modelo Zero Trust. La política bloquea todo el tráfico inicial y habilita solo comunicaciones necesarias. Con esta regla, toda comunicación no autorizada falla sin intervención del operador.

El siguiente manifiesto aplica la política a todos los Pods del namespace \texttt{production}:

\begin{lstlisting}
# default-deny-all.yaml
apiVersion: networking.k8s.io/v1
kind: NetworkPolicy
metadata:
  name: default-deny-all
  namespace: production
spec:
  podSelector: {}     # aplica a TODOS los Pods del namespace
  policyTypes:
  - Ingress           # Bloquea todo el trafico entrante
  - Egress            # Bloquea todo el trafico saliente
\end{lstlisting}

\subsubsection{Micro-segmentación por Capas de Aplicación}

% latex-paper-en · deai (Line 341) [Severity: P1] [Priority: P1]: frase genérica
Sobre Default Deny, las políticas específicas acotan el tráfico permitido. Las reglas \texttt{frontend $\rightarrow$ backend $\rightarrow$ database} usan label selectors.

En Cilium, una \texttt{CiliumNetworkPolicy} de capa 7 admite solo métodos \texttt{GET} y \texttt{HEAD} desde frontend hacia backend. En Antrea, una \texttt{ClusterNetworkPolicy} en el Tier \texttt{securityops} prioriza sobre cualquier \texttt{NetworkPolicy} de namespace.

Esta estratificación evalúa seguridad granular en condiciones reales. También cuantifica si esa granularidad introduce overhead detectable, uno de los criterios del modelo MCDA.

\subsubsection{Validación Automatizada de Políticas}

Se implementó un pipeline de validación que emplea \texttt{nc} (netcat) dentro de Pods de prueba para verificar la conectividad. Para cada par de Pods, el pipeline define una expectativa ---la conexión debe permitirse o debe bloquearse--- y reporta un error cuando el comportamiento observado no coincide con ella.

% [EDIT] Tabla booktabs y label para referencia.
\begin{table*}[htbp]
\centering
\scriptsize
\renewcommand{\arraystretch}{1.25}
\setlength{\tabcolsep}{4pt}
\caption{Casos de prueba del pipeline de validación de Network Policies.}
\label{tab:validacion-network-policies}
\begin{tabularx}{\textwidth}{|p{5cm}|Y|}
\hline
\textbf{Prueba de Validación} & \textbf{Resultado Esperado y Justificación} \\
\hline
frontend $\rightarrow$ backend (puerto 8080) & PERMITIDO --- Comunicación necesaria para el funcionamiento de la aplicación. \\
\hline
frontend $\rightarrow$ database (puerto 5432) & BLOQUEADO --- El frontend no debe acceder directamente a la base de datos. \\
\hline
backend $\rightarrow$ database (puerto 5432) & PERMITIDO --- El backend necesita leer y escribir datos. \\
\hline
database $\rightarrow$ cualquier Pod & BLOQUEADO --- La base de datos nunca debe iniciar conexiones salientes. \\
\hline
Pod externo $\rightarrow$ namespace production & BLOQUEADO --- Sin una regla explícita, ningún Pod externo puede entrar. \\
\hline
\end{tabularx}
\end{table*}

\subsection{Estructura de Repositorios y Entregables Técnicos}

% latex-paper-en · deai (Line 380) [Severity: P1] [Priority: P1]: repite contenido tabular
La Tabla~\ref{tab:repositorios-entregables} resume la organización del proyecto en tres repositorios Git.

% [EDIT] Tabla booktabs y label consistente.
\begin{table*}[htbp]
\centering
\scriptsize
\renewcommand{\arraystretch}{1.25}
\setlength{\tabcolsep}{4pt}
\caption{Estructura de los tres repositorios del proyecto y su relación con los objetivos específicos.}
\label{tab:repositorios-entregables}
\begin{tabularx}{\textwidth}{|p{3cm}|Y|p{5.2cm}|}
\hline
\textbf{Repositorio} & \textbf{Contenido Principal} & \textbf{Relación con los Objetivos} \\
\hline
\texttt{K8s-bootstrap} & Código Terraform (\texttt{00.bootstrap}), scripts de instalación K3s HA (\texttt{10.k3s-ha-do}), manifiestos de ArgoCD (\texttt{20.argo-cd-install}) y App of Apps (\texttt{30.apps-of-apps-install}). & OE1: Provee la infraestructura controlada necesaria para las pruebas. \\
\hline
\texttt{K8s-Tesis-Apps} & Overlays Kustomize para los 4 CNIs, stack de observabilidad (\texttt{apps/metrics}), benchmarks iperf3 (\texttt{cni\_test\_iperf}), cliente de latencia, Network Policies y Grafana Exporter. & OE1, OE2 y OE3: Implementa las pruebas, la seguridad y el sistema de recolección. \\
\hline
\texttt{K8S-CNI-Results} & Archivos JSON/CSV con resultados por corrida, \texttt{procesador.js} (modelo MCDA) y la SPA recomendadora (\texttt{cni-recommender-spa}). & OE3 y OE4: Implementa el modelo de comparación objetiva y el prototipo. \\
\hline
\end{tabularx}
\end{table*}

% latex-paper-en · deai (Line 404) [Severity: P2] [Priority: P2]: cadena densa
La arquitectura mantiene un flujo unidireccional. \texttt{K8s-bootstrap} crea infraestructura; \texttt{K8s-Tesis-Apps} despliega aplicaciones y genera métricas. \texttt{K8S-CNI-Results} almacena, procesa y visualiza resultados.

\subsection{Desafíos de Implementación y Soluciones Adoptadas}

% [EDIT] Tabla booktabs y label consistente.
\begin{table*}[htbp]
\centering
\scriptsize
\renewcommand{\arraystretch}{1.25}
\setlength{\tabcolsep}{4pt}
\caption{Desafíos técnicos encontrados durante la implementación y soluciones adoptadas.}
\label{tab:desafios-implementacion}
\begin{tabularx}{\textwidth}{|Y|Y|}
\hline
\textbf{Desafío Encontrado} & \textbf{Solución Técnica Implementada} \\
\hline
MTU mismatch al instalar Cilium: Pods no podían comunicarse, paquetes se fragmentaban silenciosamente en la VPC. Flannel y Calico detectan automáticamente la MTU de la interfaz subyacente en sus versiones actuales; Cilium requirió configuración explícita porque su dataplane eBPF no hereda el valor de MTU del sistema operativo sin intervención. & Se calculó la MTU efectiva: 1500 (VPC) $-$ 50 (overhead VXLAN) $=$ 1450. Se configuró el parámetro \texttt{mtu: 1450} en el overlay de Kustomize de Cilium. \\
\hline
ArgoCD en estado OutOfSync tras la rotación de CNI: Pods perdían conectividad durante el cambio. & Se implementó una secuencia de rotación segura: poner ArgoCD en modo Suspended, rotar el CNI, verificar conectividad, reactivar ArgoCD. \\
\hline
Outliers frecuentes en pruebas de madrugada: DigitalOcean mostraba throttling de CPU en horas de alta demanda del datacenter. & Se configuró una ventana de pruebas exclusiva para horario laboral. El algoritmo IQR filtra automáticamente estos outliers. \\
\hline
\texttt{procesador.js} necesitaba datos de Network Policies L7 que iperf3 no puede medir (capa 4). & Se añadió una variable binaria (soporta/no soporta) al modelo MCDA para la capacidad de Network Policies L7. Cilium recibe valor 1, los demás 0. \\
\hline
\end{tabularx}
\end{table*}

\subsection{Flujo completo de ejecución}

% latex-paper-en · deai (Line 434) [Severity: P1] [Priority: P1]: repite tabla cercana
El flujo completo sigue las etapas de la Tabla~\ref{tab:flujo-implementacion}. Terraform prepara la infraestructura; K3s y ArgoCD habilitan el plano de ejecución. Los CronJobs producen datos, \texttt{procesador.js} consolida resultados y la SPA presenta la recomendación.


\section{Pruebas y Validación}

Si el diseño describe el plan y la implementación lo ejecuta, las pruebas son el momento de la verdad: el instante en que el sistema se somete a condiciones reales y se mide si el comportamiento observado coincide con el esperado. Este capítulo documenta la metodología, los protocolos, los escenarios y los resultados de las pruebas realizadas sobre los cuatro objetivos específicos del proyecto. La documentación exhaustiva de cada protocolo garantiza que los resultados sean reproducibles por cualquier investigador que disponga del mismo testbed, y que las conclusiones sean atribuibles a los CNIs y no al procedimiento de medición.

Un principio fundamental que guió todas las pruebas fue el control de variables: en cada experimento, se modificó únicamente el CNI bajo evaluación, manteniendo constantes el hardware, el sistema operativo, la versión de Kubernetes, el tamaño de los paquetes, la duración de las pruebas y las condiciones de red. Este control es lo que permite atribuir las diferencias en los resultados al CNI y no a factores externos.

\subsection{Metodología General de Pruebas}

\subsubsection{Control de Condiciones: el Escenario de Referencia (Baseline)}

Se definió un escenario de referencia con los siguientes parámetros fijos para todos los CNIs:

\begin{table*}[htbp]
\centering
\scriptsize
\begin{tabularx}{\textwidth}{|Y|Y|}
\hline
\textbf{Parámetro Controlado} & \textbf{Valor Fijo para Todo el Experimento} \\
\hline
Tipo de instancia cloud & DigitalOcean Droplet --- Ubuntu 22.04 LTS, 2 vCPU, 4 GB RAM para nodos Worker \\
Versión de Kubernetes & K3s canal stable --- la misma versión en todos los ciclos de prueba \\
Región del datacenter & fra1 (Frankfurt) --- mismo datacenter para todos los nodos en todas las pruebas \\
Red entre nodos & VPC privada DigitalOcean --- MTU 1500 bytes, sin tráfico externo \\
Duración de cada prueba iperf3 & 300 segundos (5 minutos) por corrida \\
Número de corridas por CNI & 5 corridas independientes --- permite calcular promedios y detectar variabilidad \\
Protocolo de red bajo prueba & TCP --- el protocolo dominante en aplicaciones empresariales reales \\
Herramienta de throughput & iperf3 con salida JSON (\texttt{-J}), un solo stream TCP (\texttt{-P 1}) \\
Herramienta de latencia & Script Bash con \texttt{nc -z -w 2} --- 30 muestras por corrida \\
\hline
\end{tabularx}
\caption{Parámetros del escenario de referencia (baseline) para todas las pruebas.}
\end{table*}

\subsubsection{Automatización: de la Prueba Manual a la Prueba Orquestada}

Todas las pruebas se orquestaron mediante CronJobs de Kubernetes. La secuencia de ejecución automática en cada ciclo fue:
\begin{enumerate}
\item CronJob del cliente iperf3 se activa (minutos 0 y 30). Kubernetes crea el Pod cliente en un nodo diferente al servidor (podAntiAffinity suave, weight 100).
\item Pod cliente ejecuta iperf3 durante 300 segundos y guarda el resultado JSON.
\item CronJob del cliente de latencia ejecuta 30 mediciones de TCP connect (minutos 10 y 40) y guarda su resultado JSON.
\item Al finalizar ambas pruebas, el Resource Usage Exporter (minutos 25 y 55) consulta Prometheus para extraer el consumo de CPU y RAM del agente CNI.
\item \texttt{procesador.js} actualiza \texttt{cni-data.json} con los nuevos datos procesados.
\end{enumerate}

La orquestación automática elimina la variabilidad humana de la ejecución: cada corrida sigue exactamente la misma secuencia con los mismos intervalos, de modo que las diferencias observadas entre CNIs se deben al plugin y no al operador.

\begin{quote}
\textbf{¿Por qué 5 corridas y no 1 o 100?}
Con 1 corrida, un único evento anómalo puede dominar completamente el resultado. Con 100 corridas, el tiempo de prueba se vuelve imprácticamente largo. Con 5 corridas bien distribuidas en el tiempo, el algoritmo IQR puede identificar y descartar la corrida anómala si ocurre, y el promedio de las restantes es estadísticamente representativo. Este balance entre rigor estadístico y viabilidad operativa es estándar en la investigación de redes en entornos cloud \cite{ref7,ref8}.
\end{quote}

\subsection{Pruebas de Rendimiento de Red --- Objetivo Específico 1}

\subsubsection{Protocolo de Prueba de Throughput con iperf3}

iperf3 establece una conexión entre un proceso servidor y un proceso cliente, y el cliente envía datos lo más rápido que la red permite durante el tiempo configurado. Se usó el parámetro \texttt{-P 1} (un solo stream TCP) para representar el comportamiento de una aplicación empresarial típica, no la carga máxima teórica.

\begin{lstlisting}
# Comando iperf3 ejecutado por el CronJob cliente
# -c: servidor  -t: 300s  -J: JSON  -P 1: un solo stream
iperf3 -c iperf-server-svc.cni-benchmark.svc.cluster.local \
  -t 300 -J -P 1 \
  > /results/throughput-$(date +%Y%m%d_%H%M%S)-${CNI_NAME}.json
# El output JSON incluye:
# - bits_per_second (sender): throughput desde el emisor
# - bits_per_second (receiver): throughput real recibido
# - retransmits: paquetes TCP que debieron reenviarse
# - mean_rtt: RTT promedio en microsegundos
\end{lstlisting}

\subsubsection{Resultados de Throughput TCP por CNI}

Los resultados presentados son los promedios de las 5 corridas de 300 segundos para cada CNI, después de aplicar el filtro IQR para eliminar outliers.

\begin{table*}[htbp]
\centering
\scriptsize
\begin{tabularx}{\textwidth}{|Y|Y|Y|Y|Y|}
\hline
\textbf{Métrica} & \textbf{Flannel (VXLAN)} & \textbf{Calico (VXLAN)} & \textbf{Cilium (eBPF)} & \textbf{Antrea (OVS)} \\
\hline
Throughput Sender promedio & 9.21 Gbps & 9.48 Gbps & 9.67 Gbps & 9.35 Gbps \\
Throughput Receiver promedio & 9.18 Gbps & 9.45 Gbps & 9.64 Gbps & 9.31 Gbps \\
Retransmisiones TCP (promedio) & 12.4 / sesión & 4.8 / sesión & 2.1 / sesión & 7.3 / sesión \\
Diferencia Sender-Receiver & $\sim$0.03 Gbps & $\sim$0.03 Gbps & $\sim$0.03 Gbps & $\sim$0.04 Gbps \\
Variabilidad entre corridas ($\sigma$) & $\pm$0.18 Gbps & $\pm$0.09 Gbps & $\pm$0.07 Gbps & $\pm$0.12 Gbps \\
Outliers descartados por IQR & 1 de 5 (20\%) & 0 de 5 (0\%) & 0 de 5 (0\%) & 1 de 5 (20\%) \\
\hline
\end{tabularx}
\caption{Resultados de throughput TCP por CNI (promedio de 5 corridas de 300 segundos).}
\end{table*}

Todos los CNIs se acercan al límite teórico de la red VPC de DigitalOcean ($\sim$9.7--9.8 Gbps). Sin embargo, las diferencias en retransmisiones son significativas: Cilium con eBPF genera apenas 2.1 retransmisiones por sesión, mientras que Flannel genera 12.4, casi 6 veces más. Esta diferencia se explica porque eBPF opera directamente en el kernel y puede tomar decisiones de enrutamiento mucho más rápido que el mecanismo VXLAN de Flannel.

\subsubsection{Protocolo de Prueba de Latencia TCP}

La latencia de establecimiento TCP mide el tiempo que tarda en completarse el handshake de tres vías de TCP, que ocurre cada vez que una aplicación inicia una nueva conexión. El cliente ejecuta \texttt{nc -z -w 2} al servicio iperf3-server, mide el tiempo en ms y calcula min/avg/max/mdev sobre 30 muestras por corrida. Esta métrica es especialmente relevante para arquitecturas de microservicios, donde el costo de iniciar una conexión puede dominar el tiempo total de respuesta y el throughput del enlace no captura esa penalización.

\subsubsection{Resultados de Latencia TCP por CNI}

\begin{table*}[htbp]
\centering
\scriptsize
\begin{tabularx}{\textwidth}{|Y|Y|Y|Y|Y|}
\hline
\textbf{Métrica de Latencia} & \textbf{Flannel (VXLAN)} & \textbf{Calico (VXLAN)} & \textbf{Cilium (eBPF)} & \textbf{Antrea (OVS)} \\
\hline
Latencia mínima (ms) & 0.412 ms & 0.287 ms & 0.198 ms & 0.341 ms \\
Latencia promedio (ms) & 0.687 ms & 0.431 ms & 0.312 ms & 0.512 ms \\
Latencia máxima (ms) & 2.341 ms & 0.891 ms & 0.543 ms & 1.124 ms \\
Jitter / Desviación estándar (ms) & $\pm$0.389 ms & $\pm$0.142 ms & $\pm$0.087 ms & $\pm$0.201 ms \\
Percentil P95 (ms) & 1.834 ms & 0.712 ms & 0.478 ms & 0.893 ms \\
Percentil P99 (ms) & 2.156 ms & 0.843 ms & 0.521 ms & 1.067 ms \\
\hline
\end{tabularx}
\caption{Resultados de latencia TCP por CNI (30 muestras $\times$ 5 corridas, outliers IQR eliminados).}
\end{table*}

Los datos de latencia revelan la mayor brecha de rendimiento entre CNIs. Cilium logra una latencia promedio de 0.312 ms, mientras que Flannel alcanza 0.687 ms, el doble. El P99 de Flannel (2.156 ms) significa que el 1\% de las conexiones tardan más de 2 ms en establecerse. En una aplicación que realiza 10.000 conexiones por segundo, ese 1\% representa 100 conexiones lentas por segundo.

\subsection{Pruebas de Seguridad con Network Policies --- Objetivo Específico 2}

\subsubsection{Matriz de Pruebas de Network Policies}

Se diseñó una matriz de pruebas que cubre los escenarios críticos de seguridad en un clúster Kubernetes. La herramienta de verificación principal fue \texttt{kubectl exec} combinada con \texttt{nc} (netcat) dentro de Pods de prueba.

\begin{table*}[htbp]
\centering
\scriptsize
\begin{tabularx}{\textwidth}{|Y|Y|}
\hline
\textbf{Caso de Prueba} & \textbf{Descripción, Herramienta y Criterio de Éxito} \\
\hline
TC-SEC-01: Default Deny Ingress & Ningún Pod externo puede enviar tráfico al namespace sin regla explícita. Verificación: \texttt{nc -zv <pod-ip> 8080}. ÉXITO: timeout en menos de 5 s. \\
TC-SEC-02: Default Deny Egress & Los Pods en producción no pueden iniciar conexiones salientes no autorizadas. Verificación: \texttt{curl --max-time 3}. ÉXITO: timeout. \\
TC-SEC-03: Frontend$\rightarrow$Backend permitido & Pod con etiqueta \texttt{tier=frontend} puede conectarse al backend en puerto 8080. ÉXITO: Connection succeeded. \\
TC-SEC-04: Bloqueo Frontend$\rightarrow$Database & Pod frontend NO puede conectarse a la base de datos directamente. ÉXITO: Connection refused o timeout. \\
TC-SEC-05: Backend$\rightarrow$Database permitido & Pod backend SÍ puede conectarse a la base de datos en puerto 5432. ÉXITO: Connection succeeded. \\
TC-SEC-06: Database NO inicia conexiones & La base de datos nunca inicia conexiones salientes. ÉXITO: Bloqueo por política de Egress. \\
TC-SEC-07: Bloqueo entre namespaces & Un Pod de namespace staging no puede comunicarse con Pods de production. ÉXITO: Bloqueo confirmado. \\
TC-SEC-08: GitOps Auto-restauración & Si una política se elimina manualmente, ArgoCD la restaura en menos de 3 minutos. ÉXITO: Política restaurada automáticamente. \\
\hline
\end{tabularx}
\caption{Matriz completa de pruebas de seguridad con Network Policies.}
\end{table*}

\subsubsection{Resultados de las Pruebas de Seguridad por CNI}

\begin{table*}[htbp]
\centering
\scriptsize
\begin{tabularx}{\textwidth}{|Y|Y|Y|Y|Y|}
\hline
\textbf{Caso de Prueba} & \textbf{Flannel} & \textbf{Calico} & \textbf{Cilium} & \textbf{Antrea} \\
\hline
TC-SEC-01: Default Deny Ingress & NO SOPORTA & PASA & PASA & PASA \\
TC-SEC-02: Default Deny Egress & NO SOPORTA & PASA & PASA & PASA \\
TC-SEC-03: Frontend$\rightarrow$Backend & N/A & PASA & PASA & PASA \\
TC-SEC-04: Frontend$\rightarrow$DB bloqueado & NO BLOQUEA & PASA & PASA & PASA \\
TC-SEC-05: Backend$\rightarrow$DB permitido & N/A & PASA & PASA & PASA \\
TC-SEC-06: DB no inicia conexiones & NO SOPORTA & PASA & PASA & PASA \\
TC-SEC-07: Bloqueo entre namespaces & NO BLOQUEA & PASA & PASA & PASA \\
TC-SEC-08: Auto-restauración ArgoCD & N/A & PASA ($<$2 min) & PASA ($<$2 min) & PASA ($<$2 min) \\
Network Policies L7 (HTTP paths) & NO SOPORTA & NO SOPORTA & SOPORTA & NO SOPORTA \\
\textbf{RESULTADO GLOBAL} & \textbf{INSUFICIENTE} & \textbf{APROBADO} & \textbf{APROBADO+} & \textbf{APROBADO} \\
\hline
\end{tabularx}
\caption{Resultados de pruebas de seguridad por CNI.}
\end{table*}

El hallazgo más crítico es que Flannel no implementa Network Policies en su plano de datos. Las políticas se crean sin error en Kubernetes, pero Flannel simplemente las ignora: el tráfico que debería estar bloqueado continúa fluyendo. Esto convierte a Flannel en una opción inviable para cualquier organización con requisitos de seguridad.

Cilium sobresale como el único CNI que soporta Network Policies de capa 7 (L7). Con Cilium es posible crear reglas como ``solo permite solicitudes HTTP GET al path /api/v1/productos'', fundamental para arquitecturas de microservicios.

\subsubsection{Prueba de Overhead de Seguridad}

Se ejecutaron las pruebas de rendimiento nuevamente con las políticas de seguridad activas y se compararon con los resultados base:

\begin{table*}[htbp]
\centering
\scriptsize
\begin{tabularx}{\textwidth}{|Y|Y|}
\hline
\textbf{CNI y Estado de Políticas} & \textbf{Latencia promedio TCP / Throughput TCP} \\
\hline
Calico --- Sin políticas (baseline) & 0.431 ms / 9.45 Gbps \\
Calico --- Con políticas activas & 0.449 ms / 9.42 Gbps \\
Calico --- Overhead introducido & +0.018 ms (+4.2\%) / $-$0.03 Gbps ($-$0.3\%) \\
Cilium --- Sin políticas (baseline) & 0.312 ms / 9.64 Gbps \\
Cilium --- Con políticas activas & 0.318 ms / 9.62 Gbps \\
Cilium --- Overhead introducido & +0.006 ms (+1.9\%) / $-$0.02 Gbps ($-$0.2\%) \\
Antrea --- Sin políticas (baseline) & 0.512 ms / 9.31 Gbps \\
Antrea --- Con políticas activas & 0.538 ms / 9.27 Gbps \\
Antrea --- Overhead introducido & +0.026 ms (+5.1\%) / $-$0.04 Gbps ($-$0.4\%) \\
\hline
\end{tabularx}
\caption{Impacto de las Network Policies en el rendimiento de red de cada CNI.}
\end{table*}

El overhead de seguridad es mínimo en todos los CNIs: la latencia adicional oscila entre el 1.9\% (Cilium) y el 5.1\% (Antrea). Esto valida la premisa del diseño Zero Trust: es perfectamente viable implementar seguridad completa sin sacrificar el rendimiento de manera significativa.

\subsection{Pruebas de Eficiencia Computacional --- Consumo del Agente CNI}

Las métricas de consumo fueron capturadas por Prometheus durante las sesiones activas de iperf3 (cuando la red estaba bajo carga máxima). Se midió específicamente el consumo del proceso del agente CNI en cada nodo worker.

\begin{table*}[htbp]
\centering
\scriptsize
\begin{tabularx}{\textwidth}{|Y|Y|Y|Y|Y|}
\hline
\textbf{Métrica de Eficiencia} & \textbf{Flannel (VXLAN)} & \textbf{Calico (VXLAN)} & \textbf{Cilium (eBPF)} & \textbf{Antrea (OVS)} \\
\hline
Nombre del agente CNI & flanneld & calico-node & cilium-agent & antrea-agent \\
CPU bajo estrés (millicores) & 187 m & 143 m & 98 m & 162 m \\
CPU en reposo (millicores) & 12 m & 28 m & 31 m & 22 m \\
RAM RSS bajo estrés (MiB) & 42 MiB & 78 MiB & 185 MiB & 94 MiB \\
RAM RSS en reposo (MiB) & 31 MiB & 61 MiB & 168 MiB & 76 MiB \\
Eficiencia: Gbps por 100 millicores CPU & 4.92 Gbps & 6.61 Gbps & 9.84 Gbps & 5.77 Gbps \\
Adecuado para entornos Edge/IoT & Medio (RAM) & Medio-alto & Bajo (alta RAM) & Medio \\
\hline
\end{tabularx}
\caption{Consumo computacional del agente CNI bajo estrés de red ($\sim$9 Gbps de tráfico TCP).}
\end{table*}

Los resultados revelan un trade-off interesante. Cilium con eBPF es el más eficiente en CPU: consume solo 98 millicores para procesar el tráfico más rápido del experimento. Su métrica de eficiencia (9.84 Gbps por cada 100 millicores de CPU) es prácticamente el doble que Flannel (4.92 Gbps / 100m). Sin embargo, Cilium tiene el mayor consumo de RAM: 185 MiB bajo estrés, frente a 42 MiB de Flannel. En entornos IoT/Edge con restricciones de memoria, este consumo puede ser una limitación importante.

\subsection{Validación del Modelo Matemático MCDA --- Objetivo Específico 3}

\subsubsection{Validación del Algoritmo IQR}

Para validar que el algoritmo IQR funciona correctamente, se inyectó artificialmente 1 outlier extremo en el conjunto de datos de cada CNI (un valor de latencia 10 veces mayor al promedio real). Se verificó que el IQR lo detectara y excluyera correctamente:

\begin{table*}[htbp]
\centering
\scriptsize
\begin{tabularx}{\textwidth}{|Y|Y|}
\hline
\textbf{Verificación del IQR} & \textbf{Resultado Observado} \\
\hline
Latencia promedio de Cilium sin outlier & 0.312 ms \\
Latencia promedio de Cilium CON outlier (sin IQR) & 0.891 ms (+185\% de error) \\
Latencia promedio de Cilium CON outlier y CON IQR & 0.318 ms (+1.9\% de error) \\
¿El IQR detectó el outlier? & SÍ --- el valor 3.120 ms quedó fuera del rango [Q1-1.5$\times$IQR, Q3+1.5$\times$IQR] \\
Porcentaje de corridas descartadas en datos reales & Entre 0\% y 20\% por CNI \\
\hline
\end{tabularx}
\caption{Validación del algoritmo IQR con outlier controlado inyectado artificialmente.}
\end{table*}

Esta validación confirma que el promedio reportado para cada CNI es robusto ante eventos atípicos del entorno cloud, y que las diferencias entre CNIs en los resultados finales reflejan comportamiento estructural del plugin y no ruido experimental.

\subsubsection{Validación de la Normalización Min-Max}

\begin{table*}[htbp]
\centering
\scriptsize
\begin{tabularx}{\textwidth}{|Y|Y|Y|Y|Y|}
\hline
\textbf{Métrica (valor bruto)} & \textbf{Flannel} & \textbf{Calico} & \textbf{Cilium} & \textbf{Antrea} \\
\hline
Throughput raw (Gbps) & 9.18 & 9.45 & 9.64 & 9.31 \\
Throughput normalizado ($\uparrow$ mayor=mejor) & 0.00 & 0.59 & 1.00 & 0.28 \\
Latencia raw (ms) & 0.687 & 0.431 & 0.312 & 0.512 \\
Latencia normalizada ($\downarrow$ menor=mejor) & 0.00 & 0.67 & 1.00 & 0.46 \\
Retransmisiones raw (por sesión) & 12.4 & 4.8 & 2.1 & 7.3 \\
Retransmisiones normalizadas ($\downarrow$ menor=mejor) & 0.00 & 0.74 & 1.00 & 0.50 \\
CPU raw (millicores) & 187 & 143 & 98 & 162 \\
CPU normalizada ($\downarrow$ menor=mejor) & 0.00 & 0.49 & 1.00 & 0.28 \\
RAM raw (MiB) & 42 & 78 & 185 & 94 \\
RAM normalizada ($\downarrow$ menor=mejor) & 1.00 & 0.74 & 0.00 & 0.63 \\
Network Policies L7 (binario) & 0 & 0 & 1 & 0 \\
\hline
\end{tabularx}
\caption{Validación de la normalización min-max de las seis métricas del modelo MCDA.}
\end{table*}

La tabla de normalización muestra las seis métricas del modelo en escala comparable. Flannel es el peor en throughput, latencia, retransmisiones y CPU, pero el mejor en RAM (1.00). Cilium es el mejor en las cuatro primeras dimensiones y el único que soporta L7 (1), pero el peor en RAM (0.00). Este trade-off multidimensional es precisamente lo que justifica el enfoque MCDA con perfiles: ningún CNI domina en todas las dimensiones simultáneamente.

\subsubsection{Validación del Scoring MCDA por Perfil}

Los puntajes que se muestran a continuación son calculados por \texttt{procesador.js} a partir de los valores de precisión completa de cada corrida tras aplicar IQR, no de los promedios redondeados de las tablas de resultados. Los datos de throughput, latencia, retransmisiones y recursos de cada corrida individual se conservan en su totalidad; el redondeo presentado en las tablas anteriores puede producir diferencias apreciables al recalcular manualmente.

\begin{table*}[htbp]
\centering
\scriptsize
\begin{tabularx}{\textwidth}{|Y|Y|Y|Y|Y|}
\hline
\textbf{Perfil de Usuario} & \textbf{Flannel} & \textbf{Calico} & \textbf{Cilium} & \textbf{Antrea} \\
\hline
Fintech / Seguridad & 0.112 & 0.581 & $\star$ 0.834 & 0.542 \\
Streaming / Rendimiento & 0.298 & 0.651 & $\star$ 0.879 & 0.603 \\
IoT / Recursos limitados & 0.521 & 0.614 & 0.487 & $\star$ 0.628 \\
\hline
\end{tabularx}
\caption{Scores MCDA finales por CNI y perfil ($\star$ indica recomendación del sistema para ese perfil).}
\end{table*}

Los resultados son coherentes con las expectativas técnicas. Cilium lidera en los perfiles de seguridad y rendimiento gracias a su combinación única de eBPF y soporte de Network Policies L7. Para IoT, Antrea resulta la recomendación porque el alto consumo de RAM de Cilium (185 MiB) es un factor penalizador en entornos con memoria limitada. Este resultado demuestra que el modelo MCDA captura correctamente los trade-offs entre CNIs según el contexto.

\subsection{Validación Funcional del Prototipo Recomendador --- Objetivo Específico 4}

\subsubsection{Pruebas de Corrección Funcional}

\begin{table*}[htbp]
\centering
\scriptsize
\begin{tabularx}{\textwidth}{|Y|Y|}
\hline
\textbf{Caso de Prueba Funcional} & \textbf{Resultado Esperado vs Resultado Observado} \\
\hline
PF-01: Usuario selecciona perfil Fintech/Seguridad & ESPERADO: Cilium \#1 con score 0.834, Calico (0.581). OBSERVADO: Coincide. $\checkmark$ \\
PF-02: Usuario selecciona perfil Streaming/Rendimiento & ESPERADO: Cilium \#1 con 0.879, Calico (0.651). OBSERVADO: Coincide. $\checkmark$ \\
PF-03: Usuario selecciona perfil IoT/Recursos & ESPERADO: Antrea \#1 con 0.628, Calico (0.614). OBSERVADO: Coincide. $\checkmark$ \\
PF-04: Usuario cambia de perfil sin recargar la página & ESPERADO: El ranking se actualiza en menos de 50ms. OBSERVADO: Actualización $<$50ms. $\checkmark$ \\
PF-05: Visualización de métricas detalladas & ESPERADO: Datos expandidos correctamente (throughput, latencia, CPU). OBSERVADO: Correcto. $\checkmark$ \\
PF-06: Actualización automática de datos & ESPERADO: SPA detecta nuevo \texttt{cni-data.json} en menos de 5 minutos. OBSERVADO: 4.7 min. $\checkmark$ \\
PF-07: Responsividad en pantalla móvil & ESPERADO: Layout responsivo en iOS Safari y Android Chrome. OBSERVADO: Funcional. $\checkmark$ \\
PF-08: Manejo de error si \texttt{cni-data.json} no disponible & ESPERADO: Mensaje de espera y reintento automático. OBSERVADO: Reintento cada 30 s. $\checkmark$ \\
\hline
\end{tabularx}
\caption{Resultados de pruebas de corrección funcional del prototipo SPA.}
\end{table*}

\subsubsection{Prueba de Integridad de Datos End-to-End}

Se trazó el flujo completo de un dato específico desde su origen hasta su visualización: iperf3 reporta \texttt{bits\_per\_second: 9638201234} $\rightarrow$ script guarda el JSON sin modificaciones $\rightarrow$ \texttt{procesador.js} aplica IQR, calcula promedio = 9.638 Gbps, normaliza a 1.00, calcula score MCDA = 0.879 $\rightarrow$ \texttt{cni-data.json} registra los valores $\rightarrow$ SPA muestra ``9.638 Gbps | Score: 0.879''. Los 5 valores son idénticos. Sin redondeo, truncamiento ni pérdida de precisión en ningún punto del pipeline.

\subsection{Análisis Consolidado de Resultados}

\subsubsection{Respuesta a las Preguntas de Investigación}

\begin{table*}[htbp]
\centering
\scriptsize
\begin{tabularx}{\textwidth}{|Y|Y|}
\hline
\textbf{Pregunta de Investigación} & \textbf{Respuesta basada en evidencia} \\
\hline
¿Cuáles son las diferencias en velocidad de respuesta entre los CNIs? & Cilium tiene la menor latencia promedio (0.312 ms), un 55\% menor que Flannel (0.687 ms). En throughput, la brecha es menor (Cilium 9.64 Gbps vs Flannel 9.18 Gbps, diferencia del 5\%). La mayor diferencia está en estabilidad: el jitter de Cilium ($\pm$0.087 ms) es 4.5 veces menor que el de Flannel ($\pm$0.389 ms). \\
¿Cómo impacta cada CNI en el consumo de recursos? & Cilium es el más eficiente en CPU (98 m bajo estrés) pero el más costoso en RAM (185 MiB). Flannel es lo opuesto: ineficiente en CPU (187 m) pero liviano en RAM (42 MiB). Para entornos con poca RAM (IoT/Edge), Antrea ofrece el mejor balance. \\
¿Qué tan efectivo es cada CNI para Network Policies? & Flannel no implementa Network Policies en su plano de datos. Calico, Cilium y Antrea las implementan correctamente. Solo Cilium soporta Network Policies L7, la única opción para arquitecturas Zero Trust completas. \\
¿En qué medida se puede reducir el sobredimensionamiento? & La diferencia de 89 millicores de CPU entre Flannel y Cilium, a escala de 50 nodos, equivale a liberar $\sim$4.45 núcleos de CPU. Al costo de USD\,\$24/mes por Droplet de 2\,vCPU/4\,GB del tipo utilizado en el proyecto, liberar $\sim$2 instancias equivale a aproximadamente USD\,\$640 anuales en infraestructura DigitalOcean. \\
\hline
\end{tabularx}
\caption{Respuestas a las preguntas de investigación basadas en evidencia de pruebas.}
\end{table*}

\subsubsection{Lecciones Aprendidas de las Pruebas}

\textbf{El horario de pruebas importa.} Las pruebas ejecutadas en horario de madrugada mostraban outliers con una frecuencia 3 veces mayor que las ejecutadas en horario laboral, evidenciando mantenimientos del proveedor cloud. Sin el filtro IQR, estos outliers habrían sesgado los resultados significativamente.

\textbf{La diferencia entre promedio y percentil P99 es crucial.} El promedio de latencia de Flannel (0.687 ms) parece aceptable, pero su P99 (2.156 ms) revela que el 1\% de las conexiones experimentan latencias 10 veces mayores. Para aplicaciones donde la experiencia del usuario depende de cada conexión, este dato cambia completamente la evaluación.

\textbf{Flannel es más peligroso de lo que parece.} Es el CNI más popular en tutoriales de Kubernetes, pero las pruebas revelan que no implementa Network Policies. Muchas organizaciones están usando Flannel creyendo que tienen seguridad de red implementada, cuando en realidad sus políticas están siendo ignoradas.

\textbf{El costo de seguridad es bajo.} Las pruebas demostraron que activar Network Policies introduce apenas entre 1.9\% y 5.1\% de overhead en latencia. La diferencia es de milésimas de milisegundo. Este hallazgo refuta la creencia de que implementar seguridad de red completa requiere sacrificar rendimiento, respaldando la adopción de Zero Trust como práctica operativa estándar sin penalizar la experiencia de las aplicaciones.

\subsection{Tabla Resumen de Cumplimiento de Objetivos}

\begin{table*}[htbp]
\centering
\scriptsize
\begin{tabularx}{\textwidth}{|Y|Y|}
\hline
\textbf{Objetivo Específico} & \textbf{Estado de Cumplimiento y Evidencia} \\
\hline
OE1: Evaluar cuantitativamente el rendimiento de los 4 CNIs mediante pruebas controladas. & CUMPLIDO --- 5 corridas $\times$ 300 segundos por CNI. Datos de throughput, latencia con P95/P99 y consumo de recursos. Outliers filtrados por IQR. \\
OE2: Evaluar cómo la configuración segura de las Network Policies reduce la superficie de ataque. & CUMPLIDO --- 8 casos de prueba de seguridad ejecutados en los 4 CNIs. Hallazgo crítico: Flannel no implementa Network Policies. Cilium es el único con soporte L7. Overhead de seguridad medido: 1.9--5.1\%. \\
OE3: Desarrollar criterios y métricas para la comparación objetiva entre CNIs. & CUMPLIDO --- Algoritmo IQR validado con outlier controlado. Normalización min-max verificada. Scoring MCDA calculado para 3 perfiles. \\
OE4: Implementar un prototipo funcional que recomiende el CNI más adecuado. & CUMPLIDO --- 8 casos de prueba funcionales, todos aprobados. Prueba de integridad de datos end-to-end completada. Actualización automática validada ($<$5 min). \\
\hline
\end{tabularx}
\caption{Cumplimiento de objetivos específicos del proyecto basado en evidencia de pruebas.}
\end{table*}

\section{Resultados y Discusión}

% TODO: Este capítulo sintetiza y discute los resultados obtenidos en el capítulo anterior.
% Los datos experimentales ya están documentados en el Capítulo 4 (Pruebas y Validación).
% Aquí se desarrollará la discusión académica, la comparación con la literatura y la
% interpretación de hallazgos en el contexto del problema de investigación.

\subsection{Síntesis de Resultados Experimentales}

% TODO: Sintetizar los resultados de las Tablas 5.2-5.13 en una narrativa coherente.

\subsection{Comparación con la Literatura}

% TODO: Comparar los hallazgos con los estudios de Kang et al. (2021) \cite{ref7} y
% Koukis et al. (2024) \cite{ref8} sobre rendimiento de CNIs en entornos edge y nube.

\subsection{Discusión del Modelo MCDA}

% TODO: Discutir las implicaciones del modelo MCDA, sus limitaciones y sus aportes
% frente a métodos de comparación existentes en la literatura.

\subsection{Impacto en la Reducción del Sobredimensionamiento}

% TODO: Cuantificar el impacto económico y operativo de seleccionar el CNI adecuado
% según el perfil, con base en los datos de consumo de recursos obtenidos.

\section{Conclusiones y Trabajos Futuros}

% TODO: Las conclusiones deben sintetizar los hallazgos del proyecto en relación
% con los objetivos específicos planteados en el anteproyecto.

\subsection{Conclusiones}

% TODO: Desarrollar conclusiones para cada objetivo específico:
% OE1: Rendimiento comparativo de los cuatro CNIs
% OE2: Efectividad de las Network Policies para seguridad Zero Trust
% OE3: Validez y utilidad del modelo MCDA como herramienta de decisión
% OE4: Funcionalidad del prototipo recomendador SPA

\subsection{Trabajos Futuros}

El presente proyecto sienta las bases de un framework de evaluación que tiene múltiples posibilidades de extensión y mejora.

\subsubsection{Extensión del Framework a Entornos Edge y 5G}

El diseño actual del sistema de evaluación se implementó sobre infraestructura de nube pública. Una extensión natural sería adaptar el framework para entornos de computación edge, donde los nodos están geográficamente distribuidos con enlaces de menor capacidad y mayor latencia. Con el avance de las redes 5G y la proliferación de casos de uso de baja latencia (vehículos autónomos, cirugía remota, manufactura inteligente), la evaluación de CNIs en condiciones de edge computing se convierte en una necesidad urgente \cite{ref1}.

\subsubsection{Incorporación de Inteligencia Artificial para Selección Dinámica de CNI}

El modelo MCDA actual es estático: los pesos se definen manualmente según el perfil de usuario. Un trabajo futuro de alta relevancia sería desarrollar un modelo de aprendizaje automático que ajuste los pesos dinámicamente en función del comportamiento observado de las aplicaciones en tiempo real \cite{ref2}. Este tipo de gestión autónoma es uno de los pilares de las redes \emph{intent-based networking}.

\subsubsection{Evaluación de Nuevas Tecnologías CNI y Service Meshes}

El panorama de tecnologías de red para Kubernetes evoluciona rápidamente. Tecnologías emergentes como Kube-OVN (basado en Open Virtual Network), Multus CNI (que permite múltiples interfaces de red por Pod) y la integración de CNIs con Service Meshes como Istio o Linkerd están ganando relevancia. Un trabajo futuro significativo sería extender el framework para incluir estas tecnologías y evaluar el impacto de capas adicionales (Service Mesh) sobre las métricas de QoS.

\subsubsection{Estandarización de la Metodología de Evaluación}

En la revisión de la literatura, se identificó que no existe un estándar unificado para la evaluación de CNIs en Kubernetes. Diferentes estudios utilizan diferentes herramientas, métricas y entornos de prueba, lo que dificulta la comparación de resultados entre ellos \cite{ref7,ref8}. Un trabajo futuro de alto impacto sería publicar y proponer una metodología estandarizada basada en el framework desarrollado en este proyecto, potencialmente sometiéndola a revisión por la Cloud Native Computing Foundation (CNCF).

\subsubsection{Evaluación del Impacto Energético y Sostenibilidad}

El consumo energético de la infraestructura tecnológica es un tema de creciente relevancia económica y ambiental. Si bien este proyecto mide el consumo de CPU como proxy del consumo energético, una extensión futura sería implementar medición directa del consumo de energía mediante herramientas como PowerAPI o RAPL (Running Average Power Limit) de Intel, para calcular la eficiencia energética de cada CNI en términos de joules por gigabyte transferido.

\subsubsection{Integración con Sistemas SIEM para Análisis de Seguridad}

Las Network Policies implementadas en este proyecto proveen control de acceso a nivel de red, pero no generan alertas automáticas cuando se detectan intentos de conexión bloqueados. Un trabajo futuro sería integrar el sistema con un SIEM (Security Information and Event Management), como Elastic SIEM o Splunk, para correlacionar eventos de seguridad de red con otros indicadores de compromiso en el clúster, convirtiendo el sistema en uno activo de detección y respuesta (Network Detection and Response, NDR) \cite{ref27}.


% ============================================================
% REFERENCIAS
% ============================================================
\begin{thebibliography}{1}

\bibitem[1]{ref1}
``EndpointSlices,'' Kubernetes. [En línea]. Disponible en: \url{https://kubernetes.io/docs/concepts/services-networking/endpoint-slices/}. [Consultado: 21-ago-2025].

\bibitem[2]{ref2}
``Cluster networking,'' Kubernetes. [En línea]. Disponible en: \url{https://kubernetes.io/docs/concepts/cluster-administration/networking/}. [Consultado: 21-ago-2025].

\bibitem[3]{ref3}
``Virtual IPs and Service proxies,'' Kubernetes. [En línea]. Disponible en: \url{https://kubernetes.io/docs/reference/networking/virtual-ips/}. [Consultado: 21-ago-2025].

\bibitem[4]{ref4}
``Network plugins,'' Kubernetes. [En línea]. Disponible en: \url{https://kubernetes.io/docs/concepts/extend-kubernetes/compute-storage-net/network-plugins/}. [Consultado: 21-ago-2025].

\bibitem[5]{ref5}
D. E. Eisenbud \emph{et al.}, ``Maglev: A fast and reliable software network load balancer,'' \emph{USENIX NSDI}, 2016.

\bibitem[6]{ref6}
F. Gomes, P. Rego, y F. Trinta, ``A systematic mapping study on observability of microservices-based applications: fundamentals, classifications, and challenges,'' \emph{Computing}, vol. 107, núm. 9, 2025.

\bibitem[7]{ref7}
Z. Kang, K. An, A. Gokhale, y P. Pazandak, ``A comprehensive performance evaluation of different Kubernetes CNI plugins for edge-based and containerized publish/subscribe applications,'' Vanderbilt University, 2021.

\bibitem[8]{ref8}
G. Koukis, S. Skaperas, I. A. Kapetanidou, L. Mamatas, y V. Tsaoussidis, ``Performance evaluation of Kubernetes networking approaches across constraint edge environments,'' \emph{arXiv} [cs.NI], 2024.

\bibitem[9]{ref9}
MDPI. [En línea]. Disponible en: \url{https://www.mdpi.com/2079-9292/13/19/3972}. [Consultado: 21-ago-2025].

\bibitem[10]{ref10}
``Elevating Kubernetes network security through Cilium deployment,'' FH Joanneum. [En línea]. Disponible en: \url{https://epub.fh-joanneum.at/obvfhjhs/download/pdf/11499653}. [Consultado: 21-ago-2025].

\bibitem[11]{ref11}
W. Tu, Y.-H. Wei, G. Antichi, y B. Pfaff, ``Revisiting the Open vSwitch dataplane ten years later,'' en \emph{Proceedings of the ACM SIGCOMM 2021 Conference}, 2021.

\bibitem[12]{ref12}
``Configure MTU to maximize network performance,'' Tigera (Calico). [En línea]. Disponible en: \url{https://docs.tigera.io/calico/latest/networking/configuring/mtu}. [Consultado: 21-ago-2025].

\bibitem[13]{ref13}
``Service,'' Kubernetes. [En línea]. Disponible en: \url{https://kubernetes.io/docs/concepts/services-networking/service/}. [Consultado: 21-ago-2025].

\bibitem[14]{ref14}
Polito.it. [En línea]. Disponible en: \url{https://webthesis.biblio.polito.it/28635/1/tesi.pdf}. [Consultado: 21-ago-2025].

\bibitem[15]{ref15}
``Kubernetes networking: Comparative insights into API gateways and service mesh implementations,'' Aalto University, 2024.

\bibitem[16]{ref16}
M. Kotenko, D. Moskalyk, V. Kovach, y V. Osadchyi, ``Navigating the challenges and best practices in securing microservices architecture,'' \emph{CEUR-WS}, 2024.

\bibitem[17]{ref17}
``Roles and personas --- Kubernetes Gateway API,'' Kubernetes SIGs. [En línea]. Disponible en: \url{https://gateway-api.sigs.k8s.io/concepts/roles-and-personas/}. [Consultado: 21-ago-2025].

\bibitem[18]{ref18}
``CNI,'' \emph{cni.dev}. [En línea]. Disponible en: \url{https://www.cni.dev/docs/spec/}. [Consultado: 21-ago-2025].

\bibitem[19]{ref19}
IEEE, ``IEEE 802.1Q.'' [En línea]. Disponible en: \url{https://standards.ieee.org/ieee/802.1Q/6844/}. [Consultado: 21-ago-2025].

\bibitem[20]{ref20}
``host-local IP address management plugin,'' \emph{cni.dev}. [En línea]. Disponible en: \url{https://www.cni.dev/plugins/current/ipam/host-local/}. [Consultado: 21-ago-2025].

\bibitem[21]{ref21}
``Port-mapping plugin,'' \emph{cni.dev}. [En línea]. Disponible en: \url{https://www.cni.dev/plugins/current/meta/portmap/}. [Consultado: 21-ago-2025].

\bibitem[22]{ref22}
B. Pfaff \emph{et al.}, ``The design and implementation of Open vSwitch,'' \emph{USENIX NSDI}, 2015.

\bibitem[23]{ref23}
``Antrea --- Network Flow Visibility,'' Antrea.io. [En línea]. Disponible en: \url{https://antrea.io/docs/v1.0.0/docs/network-flow-visibility/}. [Consultado: 21-ago-2025].

\bibitem[24]{ref24}
S. Rose, O. Borchert, S. Mitchell, y S. Connelly, ``Zero Trust Architecture,'' NIST Special Publication 800-207, 2020.

\bibitem[25]{ref25}
O. Borchert, G. Howell, A. Kerman, S. Rose, y M. Souppaya, ``Implementing a Zero Trust Architecture: High-level document,'' NIST, Gaithersburg, MD, 2025.

\bibitem[26]{ref26}
R. Chandramouli y Z. Butcher, ``Building secure microservices-based applications using service-mesh architecture,'' NIST, Gaithersburg, MD, 2020.

\bibitem[27]{ref27}
J. Zhang, P. Chen, Z. He, H. Chen, y X. Li, ``Real-time intrusion detection and prevention with neural network in kernel using eBPF,'' en \emph{Proc. 54th Annual IEEE/IFIP Int. Conf. on Dependable Systems and Networks (DSN)}, 2024, pp. 416--428.

\bibitem[28]{ref28}
``CNI Essentials: Kubernetes Networking under the Hood,'' Tetrate. [En línea]. Disponible en: \url{https://tetrate.io/blog/kubernetes-networking}. [Consultado: 02-sep-2025].

\bibitem[29]{ref29}
D. E. Sierra Guerrero y H. A. Alba Castro, ``Optimización Integral de Redes Kubernetes: Comparativa de CNIs, Automatización de NetworkPolicies y Reducción de Sobredimensionamientos,'' anteproyecto de trabajo de grado, Universidad Distrital Francisco José de Caldas, Facultad Tecnológica, Ingeniería Telemática, Bogotá, Colombia, 2025.

\end{thebibliography}


% ============================================================
% BIOGRAFÍAS
% ============================================================
\begin{IEEEbiographynophoto}{Holman Alba}
Estudiante de Ingeniería Telemática.
\end{IEEEbiographynophoto}

\begin{IEEEbiographynophoto}{Duwan Sierra}
Estudiante de Ingeniería Telemática.
\end{IEEEbiographynophoto}

\end{document}
